\chapter{视觉与多模态模型}
\label{chap:vision-multimodal-models}

同一张街景照片，可以用来寻找相似图片、回答“右侧停着什么车”，也可以作为生成另一种天气下街景的条件。若再给一段录音，系统还能比较画面与声音是否相关，或把回答说出来。这些功能表面上都接受多种输入，内部需要的计算却不同。一个能比较图文相似度的模型，未必有生成句子的输出层；一个会说话的视觉助手，也未必能生成视频。

概述已经区分观测、查询和创作条件，本章在这些对象上建立后续任务共用的模型接口。我们追踪观测怎样被编码、不同来源的信息在哪里相遇、输出又怎样恢复成分数、语言或媒体。编码和连接决定模型能利用哪些证据，训练目标决定哪些对应关系得到学习，输出模块决定结果能够以什么形式出现。理解和生成可以共享其中一部分，也可以由几个专用模型协作完成。

以下沿用同一个教学场景：街景左侧有红色自行车，右侧有红色汽车，录音中可能出现发动机声。图片只提供静态外观，声音是否来自右车还需额外依据。我们可以要求选择车辆类别、回答车窗状态，也可以要求创建保留两车位置的雨天变体。这样，表示匹配、语言响应和媒体生成能够在相同对象上比较，而不会把生成的内容当作实拍证据。本章所用的小向量和概率均为教学构造，不是所介绍模型的测量结果。

\section{多模态模型的组成与接口}
\label{sec:vision-models-interfaces}

连接图像和文字之前，需要区分原始观测、网络内部表示和最终结果。它们的数值形式可能相似，却承担不同职责。例如，一张照片可以形成许多空间位置上的特征，也可以被汇总为一个整图向量；后者便于匹配，却未必足以恢复前者包含的小字和边界。

\subsection{各模态的编码与组织}
\label{sec:vision-models-encoding}

沿用前文记号，图像为$\bm{X}\in\mathbb{R}^{h\times w\times c}$，视频由带时间戳的帧$\bm{X}_t$组成。图像编码器把像素转换为$\bm{F}\in\mathbb{R}^{n_{\mathrm v}\times d_{\mathrm v}}$，每一行表示一个位置或一个汇总单元。对图像块编码，$n_{\mathrm v}$通常随空间分辨率增加。位置编码使模型能够区分“左侧的车”和“右侧的车”，但位置编码不能补回裁剪时删除的对象。

以不重叠图像块为例，设每块边长为$p$像素，且$h,w$均可被$p$整除，则共有$n_{\mathrm v}=hw/p^2$个块，不计额外汇总位置。把一个$p\times p\times c$块按固定顺序展开为$\bm x_i\in\mathbb R^{p^2c}$，线性映射$\bm A\in\mathbb R^{d_{\mathrm v}\times p^2c}$和偏置$\bm b\in\mathbb R^{d_{\mathrm v}}$形成$\bm A\bm x_i+\bm b$，再加上该块的位置表示。网络后续层在这些单元之间交换信息，得到的特征已不等于原像素。视觉Transformer（vision transformer, ViT）沿这条图像块序列计算，注意力层让不同位置的信息交互；卷积编码器则先用局部卷积逐层计算，再从空间网格取特征。不同网络可以形成相同形状的接口，却未必保留同样内容。

文字先被切成离散词元（token），即分词器按词表规则划出的字符或子词单元，再经嵌入表转为向量。音频则可以由波形转换为时频表示，再划分成局部单元编码。时频表示的两维分别描述时间和频率，和图像的横纵坐标不同。视频还需保留帧的实际时间、帧内空间位置和所属片段；单纯把第1帧、第2帧写入序列，无法区分两帧相隔0.1秒还是10秒。

若文字词表大小为$V$，第$j$个词元编号为$a_j\in\{1,\ldots,V\}$，嵌入表$\bm B\in\mathbb R^{V\times d_{\mathrm t}}$给出的输入就是第$a_j$行。分词器规定哪些字符串对应哪些编号，不能在使用时任意换表。音频波形$s_i$还需采样率$f_{\mathrm s}$，第$i$个采样点对应$i/f_{\mathrm s}$秒；将波形分窗并计算频谱后，每列仍应记录窗口的时间范围。视频按帧编码时，可以先在各帧得到视觉网格，再沿时间组织；采用时空编码器时，单元可以同时覆盖数帧，此时其时间范围也不再等于单帧时间戳。

实际系统常把连续的视觉特征也称为“视觉词元”。这里的词元表示可进入序列模型的一个计算单元，不一定经过离散词表量化。它与文字词元的整数编号、媒体生成器使用的离散编码、原始像素值，都不能直接互换。比较词元数时，必须说明每个单元对应多少空间或时间范围。

设街景原图宽1200像素、高800像素，裁出左上角为$(400,200)$、宽高为$400\times400$的区域，再调整到$200\times200$作为局部输入。以连续边界坐标计，局部结果$(x',y')=(50,80)$对应原图
\[
  (x,y)=(400+2x',\,200+2y')=(500,360)\text{。}
\]
因此，输入记录至少应保留原图编号、裁剪边界和缩放关系。若视频只选取0、2、4秒三帧，还应保留真实时间戳，而不是让下游把三帧当成连续拍摄的相邻帧。编码前的选择已经限定了模型可能看到的证据。

同一段视频还可以抽取1、3、5秒三帧。两组输入的帧数相同，却可能包含不同车辆或动作；只有保留采样位置，后续才能知道某个答案依据哪组观测。裁剪与插值也应区分：从原始高分辨率照片重新裁取车窗，可能增加相对于缩小整图的可见细节；把已经缩小的车窗插值放大，仅改变数组尺寸，不增加原始像素证据。后面的主动观察接口使用的是前一种信息获取。

\subsection{跨模态信息的连接}
\label{sec:vision-models-connection}

\term{视觉语言模型}（vision-language model, VLM）把视觉输入与语言联系起来，可以输出可比较表示，也可以生成语言。本书采用这一较宽的用法。连接音频、深度或惯性信号的模型还需相应编码和训练依据，不能仅凭“多模态”名称确定支持范围。

一种安排是两侧独立编码，再把表示映射到共同空间计算相似度。图像向量不依赖当前文字查询，适合预先建立图库。另一种安排让文字计算读取视觉特征：语言隐状态提供查询，也就是当前希望取用信息的向量；视觉特征提供用于比较的键和承载内容的值，再按匹配权重读取。此时视觉底层特征可以缓存，但读取结果随问题和已生成文字改变。还有一种安排把多源表示放入共享主干，通过位置、模态标记和注意力范围组织交互。

把第二种安排写成一次可检查的读取，有助于明确哪里依赖当前问题。设语言隐状态为$\bm H\in\mathbb R^{m\times d_{\mathrm l}}$，视觉特征为$\bm F\in\mathbb R^{n_{\mathrm v}\times d_{\mathrm v}}$，分别投影得到$\bm Q\in\mathbb R^{m\times d_{\mathrm a}}$、$\bm K\in\mathbb R^{n_{\mathrm v}\times d_{\mathrm a}}$与$\bm V_{\mathrm a}\in\mathbb R^{n_{\mathrm v}\times d_{\mathrm o}}$。交叉注意力（cross-attention）让每个文字位置按与视觉位置的相关程度加权读取视觉值：
\begin{equation}
  \bm R=\operatorname{softmax}_{\mathrm{row}}
     \left(\frac{\bm Q\bm K^\top}{\sqrt{d_{\mathrm a}}}+\bm M\right)
     \bm V_{\mathrm a}.
  \label{eq:vision-models-cross-read}
\end{equation}
这里$\bm M\in\mathbb R^{m\times n_{\mathrm v}}$在允许读取的位置取0，禁止读取的位置取$-\infty$，每行至少要有一个可读位置。行softmax使一个文字位置的权重和为1，$\bm R$的每一行是相应视觉值的加权和。问题从“车是什么颜色”变成“车窗是否打开”时，$\bm H$及查询改变，即使$\bm F$未变，读取权重也可改变。注意力数值只表示内部加权，不自动构成对象定位的标注。

这三种安排描述计算发生在哪里，不是互斥的完整模型类别。一个系统可以先用独立表示召回候选，再让联合模型逐个检查。表示对齐指训练使不同来源的内容在比较或交互中形成有用联系，参数共享指不同输入实际使用相同的权重，连接模块则负责把形状或信息路径接起来。两套编码器即使输出同样维数，也没有因此完成语义对齐。

以图库查询为例，独立编码可以保存$M$幅图的向量，新请求只需编码一次文字并比较候选。若每幅图都要与当前问题共同通过一个大型主干，通常就要执行$M$次相关联合计算。前者能否找到关系复杂的目标、后者能否在预算内处理全部候选，都需要按任务确认；模块安排只说明这些成本和能力从何而来。

\subsection{输出表示与解码}
\label{sec:vision-models-output}

输出分数、输出文字和输出媒体需要不同的预测对象。匹配接口可以计算一个向量内积；语言接口在每一步给出词表上的分布；图像或音频接口则预测像素、波形或能够被解码的中间表示。中间表示必须与解码器的训练约定相容，任意同维向量都不能送入同一个解码器。

表\ref{tab:vision-model-interfaces}固定同一街景输入，比较不同请求所需的结果。若给一个图文匹配模型附加图像生成器，匹配模型可以提供条件，但生成能力属于完整连接后的系统。反过来，视觉问答虽然输出文字，评价的仍是回答是否依据画面作出了正确判断。

\begin{table}[htbp]
  \centering
  \small
  \begin{tabularx}{\linewidth}{@{}p{22mm}XX@{}}
    \toprule
    请求 & 模型直接预测的对象 & 恢复最终结果还需什么\\
    \midrule
    选择车辆类别 & 候选描述的匹配分数 & 给定候选集与选择规则\\
    回答停车位置 & 下一文字词元的分布 & 逐词解码与来源核验\\
    生成雨天变体 & 图像潜表示的更新量 & 生成迭代与图像解码器\\
    说出判断 & 语音编码序列 & 声学解码与波形合成\\
    \bottomrule
  \end{tabularx}
  \caption{输入相同而输出接口不同的教学比较。每行描述一种完整路径，并不意味着同一个检查点同时具有四种能力。}
  \label{tab:vision-model-interfaces}
\end{table}
\FloatBarrier

\section{多模态理解模型}
\label{sec:vision-models-understanding}

如果任务是从给定候选中选择，模型可以用可比较表示形成判断。如果任务要求自由描述或回答，则还需利用语言接口组织结果。本节先展开匹配如何得到训练，再说明这种接口的用途与边界；语言响应的具体连接方式在\ref{sec:vision-models-language-generation}节展开。

\subsection{CLIP的图文双编码与相似度判定}
\label{sec:vision-models-contrastive}

为了用文字指定视觉概念，需要让“红色汽车”这样的描述与包含相应内容的图片能够比较。对比语言图像预训练（contrastive language-image pre-training, CLIP）以真实图文配对为监督，分别编码两种输入，再把正确配对与同一批次中的其他组合区分开。其图像编码器可以采用经修改的残差卷积网络ResNet或ViT，文字编码器采用Transformer，双编码结构不要求两边网络相同\scite{radford2021}

设一批有$n$对图文$(\bm{X}_i,q_i)$。图像和文字编码器分别输出$f_{\mathrm v}(\bm{X}_i)\in\mathbb R^{d_{\mathrm v}}$与$f_{\mathrm t}(q_i)\in\mathbb R^{d_{\mathrm t}}$，线性投影矩阵$\bm{W}_{\mathrm v}\in\mathbb R^{d\times d_{\mathrm v}}$、$\bm{W}_{\mathrm t}\in\mathbb R^{d\times d_{\mathrm t}}$将它们映射到$\mathbb{R}^d$。对于非零投影结果，归一化表示为
\begin{equation}
  \bm{v}_i=\frac{\bm{W}_{\mathrm v}f_{\mathrm v}(\bm{X}_i)}
                    {\lVert\bm{W}_{\mathrm v}f_{\mathrm v}(\bm{X}_i)\rVert_2},
  \qquad
  \bm{u}_i=\frac{\bm{W}_{\mathrm t}f_{\mathrm t}(q_i)}
                    {\lVert\bm{W}_{\mathrm t}f_{\mathrm t}(q_i)\rVert_2}\text{。}
  \label{eq:vision-clip-embedding}
\end{equation}
内积$s_{i,j}=\bm{v}_i^\top\bm{u}_j$就是余弦相似度。归一化使向量长度不能单独抬高分数，模型必须改变方向来改善配对。原CLIP共同学习编码器、线性投影以及控制分数尺度的温度参数。

温度$\tau>0$把相似度转换为$\ell_{i,j}=s_{i,j}/\tau$，这些未归一化分类分数称为logit。温度越小，同样的相似度差会形成越大的logit差。固定一幅图，在本批次$n$段文字中识别其配对文本，可以使用行方向交叉熵，即为正确配对的负对数选择概率计分；固定一段文字，反过来在图像中识别配对，则使用列方向交叉熵：
\begin{equation}
  \begin{aligned}
    L_{\mathrm{v\to t}}
      &=-\frac1n\sum_{i=1}^n
        \log\frac{\exp(\ell_{i,i})}{\sum_{j=1}^n\exp(\ell_{i,j})},\\
    L_{\mathrm{t\to v}}
      &=-\frac1n\sum_{i=1}^n
        \log\frac{\exp(\ell_{i,i})}{\sum_{j=1}^n\exp(\ell_{j,i})},\\
    L_{\mathrm{CLIP}}&=\tfrac12
        \bigl(L_{\mathrm{v\to t}}+L_{\mathrm{t\to v}}\bigr)\text{。}
  \end{aligned}
  \label{eq:vision-clip-loss}
\end{equation}
对角线表示数据记录中的正确配对，其他单元在该目标中充当负例。两项损失约束相同的相似度矩阵，但分母分别沿行和列求和。若批次中另有一幅也符合该描述的图，单一配对标签会把它当负例，这就是配对监督和真实语义关系之间的差别。

交叉熵为何会改变相对得分，可以从一行的目标直接推导。令$p_{i,j}=\exp(\ell_{i,j})/\sum_k\exp(\ell_{i,k})$，第$i$行的误差为$-\ell_{i,i}+\log\sum_k\exp(\ell_{i,k})$。第一项只在$j=i$时导数为$-1$，第二项的导数为$p_{i,j}$，因此两者相加得到$p_{i,j}-\mathbb I[j=i]$。若再令$\widetilde p_{i,j}$为第$j$列沿图像方向归一化的概率，完整双向平均的梯度为
\begin{equation}
  \frac{\partial L_{\mathrm{CLIP}}}{\partial\ell_{i,j}}
   =\frac{p_{i,j}+\widetilde p_{i,j}
                 -2\mathbb I[i=j]}{2n}.
  \label{eq:vision-clip-logit-gradient}
\end{equation}
正配对概率还小于1时，对角线的导数为负；负配对的导数为正。沿负梯度改变参数，便倾向于增大前者、减小后者。这里的logit不是独立可调参数，实际更新还须通过内积、归一化和两个编码器传回权重。

归一化也参与这个过程。对投影结果$\bm a\ne\bm0$与$\bm v=\bm a/\lVert\bm a\rVert_2$，其导数矩阵为$(\bm I-\bm v\bm v^\top)/\lVert\bm a\rVert_2$。这一步把沿$\bm v$方向、只改变长度的梯度分量去掉，保留能改变方向的分量。图像编码器收到同一行所有文字比较的梯度，文字编码器收到同一列所有图像比较的梯度，所以一次真实配对同时约束了两侧。零向量处公式未定义，实现需避免或用数值稳定措施处理，不能把非零条件省掉。

取一个可复算的二维教学批次：两幅图的表示分别为$(1,0)^\top$、$(0,1)^\top$，两段配对文字的表示分别为$(0.8,0.6)^\top$、$(0.6,0.8)^\top$。四个向量均为单位向量。设两组标签为“红色汽车”和“红色自行车”，并取$\tau=0.2$，就得到
\[
  \bm{S}=\begin{bmatrix}0.8&0.6\\0.6&0.8\end{bmatrix},
  \qquad
  \bm{\ell}=\begin{bmatrix}4&3\\3&4\end{bmatrix}\text{。}
\]
每一行的正确配对概率为$e^4/(e^4+e^3)\approx0.7311$，交叉熵为$\log(1+e^{-1})\approx0.3133$；列方向在这个对称例子中相同。行损失对第一行两个logit的梯度，暂不计批次平均系数，为$(-0.2689,0.2689)$。梯度下降会提高正确配对相对错误配对的分数。图\ref{fig:vision-clip-matrix}标出了两个归一化方向。

完整双向平均中$n=2$，式\eqref{eq:vision-clip-logit-gradient}给出第一行导数$(-0.1345,0.1345)$。若只为观察分数变化，暂把logit当独立变量，并取教学步长0.1，则$(4,3)$更新为约$(4.0134,2.9866)$，正确选择概率增至约0.7363。这不是真实网络的一次训练结果，因为真实参数共用、单位向量受约束；它只复算目标希望把两个分数差推向哪个方向。

训练时，先把数据记录中的图片按模型约定裁剪、缩放和归一化，把对应文本按固定分词器编码，再抽取一批记录构造整个$\bm S$。原CLIP从随机初始化的编码器开始学习，用Adam优化器调整编码器、投影和分数尺度，并按预定训练日程停止，而不是等到全部配对概率变成1。温度通常通过学习对数尺度保持为正；它改变softmax的尖锐程度，不改变未缩放余弦分数的排序。训练后保留图像处理方式、分词器、两侧权重、投影和尺度，丢弃当前批次的负例矩阵。面对新数据的正确率仍需另行评价，下降的训练目标没有给出泛化保证。

\begin{figure}[htbp]
  \centering
  % 教学数据由四个单位向量的内积构造，温度为0.2。
\begin{tikzpicture}[x=1cm,y=1cm,font=\figurefont\small,
  arr/.style={-{Stealth[length=2mm]},BookInk,line width=0.8pt}]
  \node[anchor=west] at (0,4.5) {缩放后分数$\ell_{i,j}=s_{i,j}/\tau$};
  \node[align=center] at (3.8,3.55) {$q_1$\\红色汽车};
  \node[align=center] at (6.0,3.55) {$q_2$\\红色自行车};
  \node[anchor=east] at (2.6,2.5) {图像1};
  \node[anchor=east] at (2.6,1.2) {图像2};
  \foreach \x/\y/\val in {3.8/2.5/4,6.0/2.5/3,3.8/1.2/3,6.0/1.2/4}{
    \draw[BookRule,line width=0.5pt,fill=BookPaper]
      (\x-1,\y-0.55) rectangle (\x+1,\y+0.55);
    \node at (\x,\y) {$\val$};
  }
  \draw[BookTeal,line width=1.2pt] (2.8,1.95) rectangle (4.8,3.05);
  \draw[BookTeal,line width=1.2pt] (5.0,0.65) rectangle (7.0,1.75);
  \draw[arr] (7.2,2.5) -- (8.15,2.5);
  \node[anchor=west,align=left] at (8.2,2.5) {行softmax\\$0.7311$\\$0.2689$};
  \draw[arr] (3.8,0.4) -- (3.8,-0.2);
  \node[align=center,anchor=north] at (3.8,-0.3)
    {列softmax\\$(0.7311,\;0.2689)^\top$};
  \node[anchor=west] at (0,-1.35) {对角线：数据配对\qquad 非对角线：批次负例};
\end{tikzpicture}

  \caption{CLIP双向目标的两对图文教学示例。行固定图像、列固定文字，粗框标出训练配对。行、列softmax分别形成两个分类问题，图中的0.7311是批次内选择概率，不是“图中确有该物体”的校准概率。}
  \label{fig:vision-clip-matrix}
\end{figure}

使用时，可以把类别$c$写成描述模板$q_c$，缓存其表示$\bm{u}_c$，再对新图像选取$\argmax_c\bm{v}^\top\bm{u}_c$。写“这是一张红色汽车的照片”和只写“汽车”，提供给文字编码器的上下文不同，得到的方向也可能不同。若任务只有“汽车”“自行车”两个候选，模型一定会从中选一个，除非另加拒绝规则；加入“公交车”后，归一化分数的分母又会改变。

这种零样本分类不使用当前类别任务的标注样本去拟合分类器，但仍依赖大规模图文预训练。比较结果时应记录候选集合、模板、图像处理和模型版本。整图相似度没有直接给出车辆的框、实例掩码或自由答案，这些输出需要其他接口。对比学习用正对应与候选参照构成相对判定，候选来源和数量会改变训练信号；通用目标与采样问题可参考\BookChapterRef{chap:data-derived-knowledge-acquisition}{卷八“从数据中获得知识”章}，这里的关键是图文两侧如何共同构成配对判定。

\subsection{ImageBind的图像配对与多模态表示联系}
\label{sec:vision-models-imagebind}

图文配对之外，视频帧和同时录到的声音、彩色图与对应深度、佩戴设备的画面与惯性信号，也提供了跨来源联系。若必须为所有模态两两采集配对数据，模态种类增加后组织成本会迅速上升。ImageBind用图像作为配对中心，把多个编码器接到已有视觉语义空间中\scite{vision-imagebind2023}

原实现分别编码图像、文字、音频、深度、热成像和惯性信号。图像与视频共用视觉编码器，视频输入由短时间范围内的帧构成。音频转换为梅尔频谱（mel-spectrogram），即按接近听觉尺度的频率分组记录各时间窗的能量，再划块编码。深度预处理为表示距离关系的单通道视差图，热成像也按单通道图组织；惯性信号则把加速度计和陀螺仪的时间序列通过一维投影送入Transformer。各编码器末端有独立线性投影，输出归一化到同一维数。

以声音为例，一批训练数据包括画面$\bm{X}_i$与对应声音$a_i$。模型比较$\bm{v}_i$和$\bm{u}^{\mathrm a}_j$，沿图像到音频、音频到图像两个方向使用与式\eqref{eq:vision-clip-loss}同类的配对目标。深度、热成像和惯性数据各自与其配对画面计算目标，不要求一条记录同时拥有全部模态。原论文实现冻结已预训练的图像、文字编码器，学习音频、深度、热成像和惯性编码路径。

更具体地说，设模态$m$的数据集为$\mathcal D_m=\{(\bm X_i,a_i^{m})\}$，其编码与投影参数为$\bm\theta_m$。每次在一个已存在的配对数据集中抽取批次，计算图像向量与该模态向量的双向目标$L_m$；误差经归一化和投影反传到$\bm\theta_m$。冻结图像端使误差只推动新增模态去适配图像空间，文字端也不因音频批次改变。原论文按模态分别给出AdamW日程，音频、深度和热成像设置为64轮、惯性设置为8轮，并非每条记录同时联合全部模态。训练后保留各自的预处理、编码器与投影；推理不再需要配对标签和批次内负例。

继续使用街景例子，假设发动机录音形成单位向量$(0.9,\sqrt{0.19})^\top$，图库中两幅候选的图像表示为$(1,0)^\top$和$(0,1)^\top$。两个内积约为0.9和0.4359，排序先返回前者。若两幅候选都含车辆但只有后者是这段录音的真实来源，几何排序仍可能选错；训练配对学到的是所覆盖数据中的联系，并非每次录音的声源认证。检索应返回候选编号与分数，才能继续核查内容和来源，而不是仅输出“找到了发动机”。

固定图像和文字空间的意义在于，新增模态需要向已经形成的坐标关系适配，而不是让所有编码器同时任意移动。训练仍需真实对应和足够数据覆盖。画面中既有车又有树而录音只有风声时，整段音画配对也没有显式指出声音来自哪个区域。

使用声音检索画面时，只需编码声音，再按它与图库向量的内积排序。使用声音选择文字类别时，可以比较声音向量与候选文字向量，即使新增音频编码器没有直接用音文对训练。论文在所测数据中观察到了这种未直接配对模态之间的联系；共同空间本身不是保证联系正确的定理。

一个几何事实有助于理解这种联系的可能性。若某段声音和某段文字都与同一个图像向量很接近，即$\lVert\bm{u}^{\mathrm a}-\bm v\rVert_2\leq\epsilon_{\mathrm a}$、$\lVert\bm{u}^{\mathrm t}-\bm v\rVert_2\leq\epsilon_{\mathrm t}$，三角不等式给出
\[
  \lVert\bm{u}^{\mathrm a}-\bm{u}^{\mathrm t}\rVert_2
  \leq\epsilon_{\mathrm a}+\epsilon_{\mathrm t}\text{。}
\]
但实际训练不能让每个语义一致的样本都围绕同一个准确图像锚点，也不能消除背景、视角和声源歧义，因此这只是有条件的几何解释。ImageBind的音频到图像生成展示还外接了图像生成系统，用声音表示替代相应条件；ImageBind本体仍输出表示。

\subsection{理解结果的形成与接口选择}
\label{sec:vision-models-understanding-interface}

街景查询“找相似画面”可以直接返回向量排序。“描述发生了什么”没有固定候选句表，需要根据视觉内容预测一串文字。“指出右侧车辆”则要求把表达对应到区域，仅有一句“右边有一辆车”尚未给出可检查的位置。三者都涉及理解，但需要的结果不同。

相似度适合在已定义候选中选择，语言生成适合组织可变长度的描述和答案，定位接口则要输出框、点或掩码。语言模型也可以按约定输出坐标，但仍需训练依据、坐标转换和定位评价。具体的语言指代与区域对应见\ref{sec:vision-understanding-grounding}节。

同一表示也可服务生成。图像语义向量可以作为生成器的条件，图文相似度的梯度可以调整采样方向，生成结束后的相似度还可以筛选候选。这些用途分别改变生成输入、生成过程和选择规则，不能混为“CLIP生成图像”。下一节从预测对象和训练模块出发，解释这些不同连接。

\section{多模态生成模型}
\label{sec:vision-models-generation}

生成描述、答案和语音响应，主要需要把已观察到的内容表达出来；生成图像、视频和环境声音，还要预测完整媒体中的大量细节。两类输出可以在一个系统中结合，所需目标和解码器却不同。本节分别追踪语言响应和媒体内容的生成路径。

\subsection{语言响应的生成}
\label{sec:vision-models-language-generation}

对街景问题“右侧停着什么车”，需要从同一视觉证据形成长度可变的答案，而不是从两个类别描述中选一个。设观测经过编码得到视觉表示$\bm{F}$，请求为$q$，回答的词元为$y_1,\ldots,y_m$。条件语言生成通常按此前已经产生的前缀逐词展开：
\begin{equation}
  p_{\bm\theta}(y_1,\ldots,y_m\mid\bm F,q)
  =\prod_{j=1}^m p_{\bm\theta}(y_j\mid y_{<j},\bm F,q)\text{。}
  \label{eq:vision-vlm-language}
\end{equation}
训练时，把真实回答前缀提供给模型，以目标词元的负对数概率监督；使用时，前缀逐步由模型自己产生。式中的条件依赖只有通过实际连接才有可能成立。下面几种方法改变的重点，就是视觉信息怎样到达预测$y_j$所用的语言隐状态，以及这条路径上哪些权重被训练。

一条训练记录应包含图片、请求和参考回答，必要时再含先前轮次及其角色。令$\mathcal T_i$为第$i$条记录中需要监督的回答位置，用户问题和作为条件的图片不属于这组目标位置。因为参考序列的概率是逐词条件概率的乘积，取负对数后得到可相加的误差；对$n$条记录形成经验目标
\begin{equation}
  \widehat L_{\mathrm{text}}
   =-\frac1n\sum_{i=1}^n\sum_{j\in\mathcal T_i}
      \log p_{\bm\theta}(y_{i,j}\mid
          y_{i,<j},\bm F_i,q_i).
  \label{eq:vision-vlm-text-objective}
\end{equation}
可再按目标词元数归一化，但归一化口径会改变长回答的权重，训练记录应明确采用哪一种。若参考回答中的“汽车”在当前位置预测概率为0.2，这个词元贡献$-\log0.2\approx1.6094$；若概率增到0.5，贡献降为约0.6931。负对数概率促使模型提高标注词元的概率，不直接判断标注句子的事实真假。

使用时，语言头把当前位置隐向量映射为$V$个logit，softmax后得到词元分布。选择最大概率词元是贪心解码；按分布抽样可以产生不同措辞；同时保留若干前缀、按累计对数概率选择则是束搜索。选出的词元接回上下文，再计算下一位置，直到生成结束标记或达到长度上限，最后由分词器恢复字符串。教学上若“右侧”“是”“汽车”的条件概率分别为0.8、0.9、0.7，这条三词元路径的概率为0.504，尚须乘上结束词元概率才是完整回答概率。实际分词未必把这些词各切成一个单元，例子只展示从局部分布到完整序列的计算。

\subsubsection{BLIP-2的查询提取与视觉前缀}
\label{sec:vision-models-blip2}

图像编码器输出数百个特征，已有语言模型却没有学过怎样读取它们。BLIP-2在二者之间训练查询Transformer（querying transformer, Q-Former），即用一组可学习查询从图像特征中抽取信息的网络。查询在这里是模型参数，不是用户的问题；其输出是一组随图片改变的连续表示\scite{vision-blip22023}

设图像特征为$\bm F\in\mathbb R^{n_{\mathrm v}\times d_{\mathrm v}}$，查询矩阵为$\bm Q_0\in\mathbb R^{r\times d_{\mathrm q}}$。查询通过自注意力交换信息，并以交叉注意力读取$\bm F$，形成$\bm Z\in\mathbb R^{r\times d_{\mathrm q}}$。原论文使用32个查询，无论图像编码器输出多少特征，交给下一模块的查询输出数量保持固定。压缩减少后续语言计算所读的视觉单元，也要求查询选择并保留与任务相关的信息。

这里可以用式\eqref{eq:vision-models-cross-read}跟踪一次查询读取。把$32\times768$查询隐表示投影为查询，把冻结图像特征投影为键和值，就得到$32\times n_{\mathrm v}$的权重矩阵。每行读取图像网格的一组加权值，再通过残差与后续层形成查询输出。以两个视觉值$(1,0)^\top$、$(0,1)^\top$及单个教学查询为例，若读取权重为$(0.25,0.75)$，输出贡献就是$(0.25,0.75)^\top$；这是信息汇聚，不是自动裁出右车。Q-Former的共享自注意力层从预训练BERT初始化，新增交叉注意力随机初始化，查询本身也作为参数学习。已有图像编码器来自CLIP或EVA-CLIP，后接的语言权重来自OPT或FlanT5，不是所有模块都从零开始。

第一阶段连接冻结图像编码器和Q-Former，使用图文对同时学习三种目标。对比目标让图像查询和文字分别编码，二者在自注意力中不能互看，以免把待匹配文本直接泄漏给图像表示；每段文字与多个查询输出比较，取最高相似度作为图文分数。配对判定目标则允许查询和文字双向交互，由各查询输出的二分类logit取平均，判断图文是否匹配，并使用困难负例。

第三种目标是图像条件文字生成。查询之间可以互看，但不能读取目标文字；每个文字词元可以读取全部查询和此前文字。由于文字不能直接读取冻结图像特征，生成目标迫使Q-Former经过查询路径传递有用信息。三种目标共用Q-Former参数，但注意力掩码不同，分别约束独立匹配、细致核验和信息表达。

设归一化查询输出为$\bm z_{i,k}$、文字表示为$\bm t_j$，独立匹配的分数是$\max_{1\leq k\leq r}\bm z_{i,k}^{\top}\bm t_j$，再按式\eqref{eq:vision-clip-loss}形成批次目标。配对判定则对一条图文组合给出两类概率$(p_0,p_1)$，真实配对标签$b=1$、采样负例标签$b=0$，误差为$-b\log p_1-(1-b)\log p_0$。条件文字生成用式\eqref{eq:vision-vlm-text-objective}，但此时预测者是Q-Former的文字分支。三个可计算误差相加，再把梯度传回共享层、查询和各目标的输出层，图像编码器权重不变。最大相似度让最相关查询承担该配对的直接比较梯度，并不强制32个查询分别对应32类对象。

第二阶段用$\bm W\in\mathbb R^{d_{\mathrm l}\times d_{\mathrm q}}$把$\bm Z$线性投影到语言模型的嵌入维数$d_{\mathrm l}$，得到$\bm P=\bm Z\bm W^\top\in\mathbb R^{r\times d_{\mathrm l}}$。这些向量放在文字输入前面，称为\term{视觉前缀}（visual prefix）：它们占据语言模型的输入位置，用连续向量提供图像条件，却不对应可直接打印的文字。冻结图像编码器和大型语言模型后，生成损失训练Q-Former与投影层。对于解码器式语言模型，目标是根据视觉前缀生成图像描述；对于编码器—解码器式语言模型，论文还使用前缀语言建模，把部分文字与视觉前缀交给编码器，预测剩余文字。

第一阶段保留训练好的Q-Former作为第二阶段初始化，另加语言宽度的投影。每个批次重算图像查询、视觉前缀和参考文字的概率，误差经冻结语言层反传到$\bm P$，再更新投影与Q-Former。原论文用AdamW和预定学习率日程，第一阶段250K步、第二阶段80K步；这些停止点是训练配置，不是证明目标已全局最小。图文描述可以来自原始网页，也可以由论文的描述生成与筛选流程获得，合成文字不应被当作无误标注。最终保留图像编码器、查询模块、投影、语言模型与分词器，预测不再运行第一阶段的配对二分类头。

在街景问答中，模型先计算图像特征，再取得32个查询输出，接上“右侧停着什么车”及相应提示，由语言模型逐词生成答案。固定图片的视觉前缀可用于不同文字请求，但压缩时未保留的细节无法仅靠换一种问法恢复。这里的冻结说明针对两阶段预训练，不能推广到原论文所有下游微调设置，更不能作为所有VLM的固定训练顺序。

\subsubsection{Flamingo的视觉重采样与交错上下文}
\label{sec:vision-models-flamingo}

若输入包含“示例图像—示例回答—待答图像—问题”，视觉信息还需要与文字在序列中的位置对应。Flamingo先把每幅图像或视频压缩为固定数量的视觉单元，再让语言模型在中间层读取这些单元。原模型冻结已有视觉编码器与语言模型，训练新增连接部分\scite{vision-flamingo2022}

视觉编码器输出空间特征；视频则对采样帧分别编码，并加入时间嵌入。感知重采样器（Perceiver Resampler）使用可学习查询，通过注意力把数量可变的特征重采样为64个输出。这里的“重采样”是学习式信息汇聚，不是从原图均匀抽取64个像素。其注意力的键和值还包含查询隐表示本身，使这些潜在单元能够参与后续汇聚。

与BLIP-2把视觉表示放在语言输入前面不同，Flamingo在冻结的语言层之间插入门控交叉注意力与前馈模块。交叉注意力的查询来自当前语言隐状态，键和值来自重采样后的视觉表示。若用$C(\bm H,\bm Z)$概括一个新增交叉注意力分支，残差更新的核心形式为
\begin{equation}
  \bm H'=\bm H+\tanh(\alpha)\,C(\bm H,\bm Z)\text{，}
  \label{eq:vision-flamingo-gate}
\end{equation}
其中$\alpha$是从0开始学习的标量门控。初始时新增分支对输出的贡献为0，语言模型保留原来的计算；门控和连接权重随后从文字预测损失中学习。完整模块还包含相应的前馈分支，不应把式\eqref{eq:vision-flamingo-gate}当成整层结构。

初始贡献为0不意味着门控也学不到信号。若把后续误差对$\bm H'$的导数记为$\bm G$，则$\partial L/\partial\alpha=(1-\tanh^2\alpha)\langle\bm G,C(\bm H,\bm Z)\rangle$，其中内积对矩阵全部元素求和。$\alpha=0$时这个导数仍可非零；而该分支权重的直接梯度先乘以$\tanh\alpha$，在门控打开后才逐渐获得信号。教学上若原隐向量为$(1,0)^\top$、分支贡献为$(0,2)^\top$，门控从0变为0.2后，更新结果为$(1,0.3948)^\top$。这说明新视觉信息怎样进入冻结语言计算，不是说每个门控实际都朝正方向变化。

交错输入还需要控制可见范围。原Flamingo的每个文字位置直接交叉注意到它之前最近的一幅图像或一个视频的视觉表示，而不是同时直接读取全部历史图片。历史图像仍可通过此前文字隐状态和语言自注意力间接影响结果。因此，示例中的红车和待答图中的红车虽然类别相同，也必须保持来源位置，不能把示例属性无条件移到目标上。

训练数据包含网页中的交错图文、图文对及视频文字对，目标仍是对可见视觉上下文条件下的文字进行预测。使用时，把几个已知示例和新请求放入上下文即可，示例本身不引起梯度更新。它们帮助说明当前要做什么任务，但示例数、输入处理和上下文压缩仍限制能保存的细节。图\ref{fig:vision-language-connectors}把前缀、层间读取与直接投影三条路径放在一起，便于区分视觉信息实际进入语言计算的位置。

一条交错训练记录还要保存视觉单元插入文字的位置，才能形成上述可见掩码。对每个数据来源分别计算文字负对数概率，再按来源权重累积梯度，更新重采样器、新增门控连接层及相应新增词元；原NFNet视觉权重和Chinchilla语言权重保持冻结。训练到预定日程后保留这些模块。推理把示例图$A$、示例答案“车窗关闭”、目标图$B$和问题依次放入上下文，目标答案位置直接读取$B$的64个视觉单元，再逐词解码；$A$的答案提供格式示范，却不能代替观察$B$的车窗。闭集任务还可对每个候选答案累积条件对数概率，按相同长度处理规则选择，而不是必须自由生成一句长回答。

\subsubsection{LLaVA的特征投影与视觉指令适配}
\label{sec:vision-models-llava}

若保留图像的网格特征，可以直接把每个视觉单元映射到语言嵌入空间，再让语言主干学习怎样使用这些输入。原始LLaVA采用这一连接方式，并用视觉指令数据训练回答行为。本节依据原论文v2中的线性投影设置，后来的多层感知机连接器不属于这里的模型\scite{vision-llava2023}

CLIP视觉编码器从图片取得网格特征$\bm F$，投影矩阵把各位置的特征转为与文字嵌入同宽的$\bm P$。这些视觉单元与问题词元一起进入自回归语言模型。投影并没有把每个图像块翻译成一个汉字或单词，而是调整连续特征的坐标，使语言层可以对它们计算注意力。

设$\bm W\in\mathbb R^{d_{\mathrm l}\times d_{\mathrm v}}$，则采用行存特征时$\bm P=\bm F\bm W^\top$。网格顺序和图像输入边界帮助模型知道哪些位置属于图片；语言自注意力随后在视觉位置与文字位置之间计算，和Flamingo另设一条交叉注意力分支不同。以两个二维教学网格$(1,0)$、$(0,1)$及$\bm W=\left[\begin{smallmatrix}2&0\\0&1\end{smallmatrix}\right]$为例，投影得到$(2,0)$、$(0,1)$。投影改变坐标宽度或尺度，但没有额外读取像素，也没有建立“第一行一定是自行车”的物体语义。

训练分为特征对齐和视觉指令适配。对齐阶段用图像描述构造“请简要描述图片”一类请求，冻结视觉编码器与语言模型，只更新投影层。语言模型的误差仍需沿计算图传回投影输入，冻结权重并不等于可以截断这条反向传播路径。指令适配阶段继续冻结视觉编码器，同时更新投影与语言模型，数据则包括对话、详细描述和复杂问答。

对多轮数据，输入保留图片、用户请求和前轮回答，损失监督助手应生成的回答词元。若请求是“比较两辆车的状态”，模型依赖同一视觉网格形成两个对象的描述，并逐词预测比较句。训练示例帮助模型学会“比较”这一请求形式，能否正确读取车的位置、朝向和遮挡仍由视觉证据及训练覆盖决定。

原论文的数据构造用图片已有描述和对象框向文字教师提供视觉内容的符号表示，再生成对话、详细描述和复杂问答；教师在这一流程中没有直接看原像素。训练LLaVA时仍把真实图片交给CLIP视觉编码器，不能把教师的输入方式误写成学生只读文字。第一阶段从预训练CLIP与Vicuna加载权重，随机初始化投影，用描述监督提高回答概率；第二阶段从第一阶段投影继续，用助手位置掩码同时更新投影与Vicuna，视觉权重不变。原论文使用AdamW及有限轮次日程，最终保留更新后的语言模型、投影和固定视觉编码器。

冻结视觉端时，$\bm F$可作为常量；若误差对$\bm P$的导数为$\bm G\in\mathbb R^{n_{\mathrm v}\times d_{\mathrm l}}$，投影的梯度就是$\bm G^\top\bm F$，形状与$\bm W$相同。这给出“语言误差怎样训练连接器”的直接依据。第二阶段再允许语言权重改变，模型便同时学习读取视觉前缀和组织请求答案。使用时先得到网格投影，接入问题与角色标记，再按前面的逐词规则生成回答。预测出的“右侧汽车车窗关闭”应回查目标图；语言损失只对训练回答负责，不能保证任意新图片的细节判断正确。

原模型可以生成自然语言，但没有因此得到任意精度的坐标、实例掩码或图像生成能力。若回答写了“右车正在行驶”，单张静态图片只提供姿态和场景线索，通常不足以确定实际速度；模型必须避免把语言中常见的情境补全当成观测。

\begin{figure*}[htbp]
  \centering
  % 每行仅保留用于区分接入位置的模块，不展开完整Transformer层。
\begin{tikzpicture}[x=1cm,y=1cm,font=\figurefont\small,
  block/.style={draw=FigureModel,fill=FigureModel!10,line width=0.8pt,
    minimum height=9mm,minimum width=21mm,align=center},
  input/.style={block,draw=FigureInput,fill=FigureInput!10},
  output/.style={block,draw=FigureOutput,fill=FigureOutput!10},
  arr/.style={-{Stealth[length=2mm]},BookInk,line width=0.8pt}]
  \node[anchor=west,font=\figurefont\small\bfseries] at (0,6.7) {BLIP-2：查询压缩后作为输入前缀};
  \node[input] (bf) at (1.1,5.7) {视觉特征};
  \node[block] (bq) at (4.1,5.7) {Q-Former};
  \node[block] (bp) at (7.2,5.7) {投影后的\\视觉前缀};
  \node[block] (bl) at (10.7,5.7) {语言模型};
  \node[output] (bo) at (14.1,5.7) {回答词元};
  \draw[arr] (bf) -- (bq);
  \draw[arr] (bq) -- (bp);
  \draw[arr] (bp) -- (bl);
  \draw[arr] (bl) -- (bo);
  \node[anchor=north] at (4.1,5.15) {可学习查询读取图像};
  \node (btext) at (10.7,4.7) {文字问题};
  \draw[arr] (btext) -- (bl);

  \node[anchor=west,font=\figurefont\small\bfseries] at (0,3.9) {Flamingo：在语言层之间读取视觉表示};
  \node[input] (ff) at (1.1,2.45) {视觉特征};
  \node[block] (fr) at (4.1,2.45) {Perceiver\\Resampler};
  \node[block] (fc) at (7.2,2.45) {门控交叉\\注意力};
  \node[block] (fl) at (10.7,2.45) {后续语言层};
  \node[output] (fo) at (14.1,2.45) {回答词元};
  \draw[arr] (ff) -- (fr);
  \draw[arr] (fr) -- node[above] {$K,V$} (fc);
  \draw[arr] (fc) -- (fl);
  \draw[arr] (fl) -- (fo);
  \node (fh) at (7.2,3.35) {语言隐状态$Q$};
  \draw[arr] (fh) -- (fc);
  \node[anchor=north] at (4.1,1.9) {每个图像／视频64个单元};

  \node[anchor=west,font=\figurefont\small\bfseries] at (0,0.7) {原始LLaVA：网格特征投影后进入共享序列};
  \node[input] (lf) at (1.1,-0.25) {视觉网格};
  \node[block] (lp) at (4.1,-0.25) {线性投影};
  \node[block] (ls) at (7.2,-0.25) {视觉单元\\与问题序列};
  \node[block] (ll) at (10.7,-0.25) {语言模型};
  \node[output] (lo) at (14.1,-0.25) {回答词元};
  \draw[arr] (lf) -- (lp);
  \draw[arr] (lp) -- (ls);
  \draw[arr] (ls) -- (ll);
  \draw[arr] (ll) -- (lo);
\end{tikzpicture}

  \caption{三种视觉接入方式的结构示意，按各原论文的预训练路径绘制。BLIP-2由少量查询形成视觉前缀，Flamingo在语言层中新增读取视觉表示的交叉注意力，原始LLaVA把视觉网格线性投影后交给语言主干。图中省略各模块内部层，框旁文字指出视觉信息进入的位置；冻结范围由正文所述阶段决定。}
  \label{fig:vision-language-connectors}
\end{figure*}
\FloatBarrier

\subsubsection{Qwen3-VL的多层视觉接入与视频时间组织}
\label{sec:vision-models-qwen3vl}

多页文档和视频包含分散证据，只读取视觉编码器的最末层、再把所有帧当成无时间含义的图片，可能丢失细节与先后关系。Qwen3-VL技术报告v1在动态分辨率编码、视觉合并、跨层接入和时间表示上组织这些信息。报告区分思考版与非思考版，以下机制不意味着两类后训练版本具有相同响应行为。

视觉编码器接受动态分辨率输入，输出长度随图像大小变化的视觉特征。视觉语言合并模块使用多层感知机（multilayer perceptron, MLP），即在逐位置的线性映射之间加入非线性运算；这里用两层映射把相邻$2\times2$特征压缩为一个语言宽度的视觉单元\scite{vision-qwen3vl2025}
空间合并降低上下文长度，同时也把四个位置的信息汇集到一个位置表示中；实际细节保留需结合输入分辨率和训练检查。

多层视觉接入DeepStack进一步利用视觉编码器三个不同深度的中间特征。每组特征经独立合并模块映射后，加到语言主干前三层相应的视觉位置隐状态上。视觉信息由多个深度进入语言计算，而不是全部仅依赖初始输入。残差接入不增加额外序列位置，但仍增加了特征投影和存储等计算。

位置表示有两项不同安排。交错多模态旋转位置编码（interleaved multimodal rotary position embedding, interleaved MRoPE）按各坐标旋转查询与键的部分分量，使比较分数能够反映相对位置。它把时间、高度、宽度对应的旋转位置维度分散到不同频率范围，避免把某个坐标维度集中在单一频段。视频时间戳则以文字形式放在每个时间块之前，例如“3.0秒”，提供明确的时间读数。两者共同使用，但文字时间戳不是从MRoPE隐向量中自动打印出来的数值。

在一个教学视频中，采样块对应1.0、3.0和8.0秒。保留这些时间戳，模型才有条件判断后两个观测相隔5秒；只给顺序编号会把这种间隔信息删掉。不过，如果敲门动作发生在4.2秒且没有进入任何采样块，时间戳再准确也无法补回动作。多图输入同样要保留图号，避免把第二幅图的车窗状态归到第一幅图。

报告的预训练先只更新合并模块，冻结视觉编码器和语言主干，随后转为全部模块训练，并逐步扩展上下文窗口。后训练包含指令适配、蒸馏与强化学习等安排，思考版使用相应的推理形式数据，非思考版使用标准回答形式。使用时应选定具体检查点、格式和观察预算。报告支持的256K词元上下文是容量约定，不是对任意长文档小字和短暂视频事件的保存保证。

这些阶段的样本仍需要明确输入与目标。图像描述样本以图片为条件、描述为目标，视频问答样本以带时间戳的帧组和问题为条件、答案为目标，交错文档样本按顺序保留图片与文字。它们用文字预测误差训练，但来源、目标位置及长度不同，不能把每张图都当作一个类别标签。视觉端由SigLIP-2检查点继续训练，语言端由Qwen3初始化，合并模块将空间与层间特征接入语言计算；预热阶段误差只改变合并权重，全参数阶段再共同改变三个模块。扩大上下文改变的是一次可计算序列的允许长度，并不替代更密采帧或读取原图。

多图比较时，模型先分别编码两幅街景，以图号和输入边界标记组织视觉块，再把问题接入；答案中的“第二张图车窗打开”需要依赖第二图的实际特征。视频步骤问答则在每个帧组前接入实际时间文字，模型生成事件与时间描述时使用这些序列条件。语言头形成候选词元后，仍按式\eqref{eq:vision-vlm-language}解码，时间字符串不是额外测量器自动返回的真值。最终保留具体思考或非思考检查点及其处理器；后训练产生的响应形式不能反推输入采样一定保留了所问动作。

\subsubsection{Qwen2.5-Omni的音画输入与文字语音响应}
\label{sec:vision-models-qwen-omni}

视频中的敲门动作和敲击声可能互相补充，也可能错位或来自画外。Qwen2.5-Omni技术报告v1通过Thinker--Talker结构处理这类输入并输出文字与语音。Thinker读取多模态上下文、生成文字及隐表示，Talker利用这些信息生成语音编码，再恢复成声音\scite{vision-qwenomni2025}

输入音频被转换为梅尔频谱，经音频编码器形成时间特征；图像和视频由视觉编码器处理。时间对齐多模态旋转位置编码（time-aligned multimodal rotary position embedding, TMRoPE）给音视频提供共同时间尺度，报告中一个时间位置单位对应40毫秒。音频单元沿时间排列，视频单元同时保留帧时间和空间位置。输入还按实际时间划分为2秒块，每块内先放视觉表示、后放音频表示，再按块交错。

因此，序列中相邻的视觉和声音单元不一定只相差一个物理采样点，实际时间由位置记录共同表达。这个安排也不证明两段内容来自同一个声源。若门外另有人敲击金属，画面中的手部动作与声音可能在时间上接近，却缺乏确定来源。回答应分别依据看见的动作、听见的声音和额外文字线索。

Thinker按式\eqref{eq:vision-vlm-language}的同类方式预测回答文字。Talker接收Thinker的高层隐表示以及已采样文字词元的嵌入，以双轨自回归方式组织文字和语音编码。隐表示提供语义、语气等上下文，离散文字则帮助确定具体词语，避免语义接近但读音不同的内容仅由连续表示区分。

语音编码还不是可以播放的波形。报告的流式解码路径用流匹配模型将编码转为梅尔频谱，再由声码器合成波形；\ref{sec:vision-models-flow-matching}节将解释所需的速度预测目标。解码采用有限块上下文，包含回看和一定前看，因此流式表示分块产生结果，并不意味着没有等待。

训练职责也有区别。多模态预训练起步时固定语言模型，先适配连接再训练相关编码器，随后展开全参数与长上下文训练；Thinker再用多模态指令数据适配回答。Talker单独经历语音续接、偏好优化和说话人适配等训练安排，以学习从语义上下文到语音的生成。理解监督不能替代语音发音监督，音画时间对齐也不能替代输出声学解码。

以文字与对应朗读音频为训练记录，先用语音编码器把真实声音转为离散编码，再以此前编码、文字与Thinker隐表示预测下一个语音编码；目标编码的负对数概率提供Talker的监督。它与回答文字的目标可以相互联系，但真实语音决定音色、节奏与发音的参考，不能只用一条文字回答推出全部声学目标。波形解码器另以编码和真实频谱为配对学习频谱恢复，声码器学习从频谱到波形。推理时Talker使用自己此前生成的编码，流式解码只读取约定的有限块，最终得到可播放的采样序列，而不是把Thinker的隐藏向量直接保存成音频文件。

在敲门例子中，输入路径保留动作和声音，Thinker形成“画面中的人抬手敲门，并伴有敲击声”这样的候选回答，Talker把确定的回答组织为语音。该版本支持文字与语音输出，不因此具有图像或视频生成模块。首段声音的到达时间还包含输入编码、文字生成、语音编码预测及波形解码等开销。

\subsubsection{观察决策、视觉工具返回与后续语言计算}
\label{sec:vision-models-observation}
\label{sec:vision-models-countercredit}

如果整图缩小后右车的车窗只剩几像素，语言模型继续生成更长解释仍没有得到新的观测。另一种计算是暂停回答，指定原图中的车窗区域，由工具返回该区域的像素，再继续语言预测。观察动作（visual observation action）指这次取回视觉内容的决定，不是答案中的一句“我再仔细看看”。它改变后续计算能够读取的信息，而仅延长原上下文中的文字生成不改变像素来源。

一次调用可以记录为$(h_i,c_i,o_i)$：$h_i$是调用前已有的完整状态，包括请求、先前图片与返回记录、已生成文字；$c_i$是此次动作及参数；$o_i$是实际返回图像。裁剪参数还需原始图片或先前观察的来源编号、坐标范围和输出尺寸。若在先前裁剪中再次裁剪，应组合两次坐标变换才能回到原图，而不能把第二次的局部坐标直接当作原图坐标。工具实际取出像素后，模型重新编码$o_i$，把其表示放入相应图像轮次，随后才预测后续词元。越界、空区域或调用格式无效时，没有可用的新观察，系统应保留失败记录，而不是默默把失败当成已经看到了对象。

有了这个接口，还要学习什么时候值得调用。最终答对只能说明整条路径的结果，不能说明这次裁剪有必要，也不能说明答案真的依赖返回像素。2026年预印本CounterCredit在已经发生的调用前状态上，分别检查“继续这条视觉分支比立即回答增加了多少正确答案得分”和“真实像素比同格式随机返回增加了多少得分”。本节依据其v1方法、训练配置和返回格式，范围限于返回单幅裁剪图像的工具\scite{vision-countercredit-2026}。

训练记录包含原图$\bm X$、问题$q$和参考答案$y$。策略$\pi_{\bm\theta}$从已有视觉语言检查点生成交错文字与裁剪的轨迹$\xi$，最后按任务规则取得正确性回报$r_{\mathrm{out}}(\xi)\in\{0,1\}$。为了比较同一调用的不同返回，论文不重新生成一段不同的调用前解释，而是固定已实现的$h_i$。真实分支接入原动作$c_i$与原返回$o_i$；立即回答分支从$h_i$直接要求回答；随机分支固定$h_i$与$c_i$及其文字参数，只把返回像素替换成从原图均匀采样的同尺寸图块$\widetilde o_{i,k}$。随机图块不事先筛成“保证无关”，因而可能偶然含有答案证据，这种随机参照的限制必须保留。

先明确比较用的可计算量。对于多选问题，设$O$为题目给出的选项字母集合，$e$为固定的立即作答提示，则参考答案得分为
\begin{equation}
  U_{\bm\theta}(p)
   =\frac{\pi_{\bm\theta}(y\mid p,e)}
          {\sum_{a\in O}\pi_{\bm\theta}(a\mid p,e)}.
  \label{eq:vision-countercredit-score}
\end{equation}
$p$是待探测的状态前缀，分母只在选项字母间归一化，三个分支使用相同提示与答案起始格式。得分不是最终抽样答对的频率，也不是外部判断器给整段推理打分。自由回答没有这个有限选项分母，论文改用参考答案在真实前缀监督下的每词元平均对数概率，单位为每词元奈特；它与$[0,1]$概率差的尺度不同，论文分别校准多选死区0.05与自由回答死区0.44奈特每词元。探测得分停止梯度，只用来计算策略训练信号。

令$U_i^{\mathrm{real}}=U_{\bm\theta}(h_i,c_i,o_i)$、$U_i^{\mathrm{now}}=U_{\bm\theta}(h_i)$，并用$K$个随机返回估计参照均值。两项差值为
\begin{equation}
  \begin{aligned}
    D_i&=U_i^{\mathrm{real}}-U_i^{\mathrm{now}},\\
    \overline U_i^{\mathrm{rand}}
      &=\frac1K\sum_{k=1}^{K}
          U_{\bm\theta}(h_i,c_i,\widetilde o_{i,k}),\\
    E_i&=U_i^{\mathrm{real}}-\overline U_i^{\mathrm{rand}}.
  \end{aligned}
  \label{eq:vision-countercredit-values}
\end{equation}
$D_i$称为决策价值（decision value），比较实际视觉分支与立即作答；$E_i$称为证据价值（evidence value），比较实际像素与同动作下的随机返回。前者还包含动作格式或延后回答的影响，后者尽量保持这些因素不变。它们检查的是当前策略的参考答案得分变化；即使原来已经能答对，增加置信度仍可能产生正值，因此论文的“需要”不能直接改写成逻辑上“不看必错”。

只奖励其中一项会遗漏不同失败情形。若$D_i$大而$E_i$接近0，随机图块也带来同样收益，这次调用没有显示真实像素的额外作用。若$E_i$大而$D_i$接近0，像素相对于随机返回有信息，但立即作答已同样充分。论文因此取$a_i$为探测资格指示，包括有效单图返回、存在不同随机图块和得分有限，取$\varepsilon$为预先校准的误差死区，并定义
\begin{equation}
  \begin{aligned}
    b_i&=\min\{E_i,D_i\}
        =U_i^{\mathrm{real}}
          -\max\{U_i^{\mathrm{now}},\overline U_i^{\mathrm{rand}}\},\\
    v_i&=r_{\mathrm{out}}(\xi)\,a_i\,\mathbb I[b_i>\varepsilon].
  \end{aligned}
  \label{eq:vision-countercredit-verification}
\end{equation}
这里$v_i=1$才表示此次调用获得双重核验。第二个等式来自同时减去两个参照时取较小差值，等价于比较较高参照得分；所以任何一边没有超过死区都不能通过。正确性因子还要求整条轨迹最终答对。$\varepsilon$用训练前的随机图块对照校准后固定，避免把有限$K$导致的小波动全当成收益，且不同得分单位使用不同值。

仍以车窗问题作教学计算。假设立即回答的正确选项得分为0.40，实际裁剪后为0.80，三个随机返回的得分为0.35、0.45、0.40，则$\overline U^{\mathrm{rand}}=0.40$，$D=E=b=0.40$。取多选配置$\varepsilon=0.05$，有效调用且最终答对时即可通过。若实际与随机返回都为0.80，则$D=0.40$而$E=0$，不能通过；若立即回答为0.82、实际为0.80、随机均值为0.40，则$E=0.40$而$D=-0.02$，也不能通过。三种情形固定了同一个问题及像素来源，展示两个信号为何不能互相替代。

调用收费还要防止模型仅靠减少成本而偏好错误答案。设$n(\xi)$为已执行调用数，$\mathcal V(\xi)=\{i:v_i=1\}$，核验通过的调用获得超过门槛部分的返还，未通过的调用付费。论文的轨迹价格与总回报为
\begin{equation}
  \begin{aligned}
    P(\xi)
      &=\min\left\{C,\,
         \gamma\sum_{i\in\mathcal V(\xi)}(b_i-\varepsilon)\right\}\\
      &\quad-\rho\bigl(n(\xi)-|\mathcal V(\xi)|\bigr),\\
    R(\xi)&=r_{\mathrm{out}}(\xi)+P(\xi).
  \end{aligned}
  \label{eq:vision-countercredit-price}
\end{equation}
$\gamma$把得分收益换成回报，$C$是每条轨迹返还上限，$\rho>0$是每个未核验调用的费用。失败调用和错误答案轨迹中的调用也计费，不通过便不能免费。若每条轨迹最多$n_{\max}$次调用且$\rho n_{\max}<1$，正确轨迹至少获得$1-\rho n_{\max}>0$，错误轨迹没有返还、总回报至多0，因此费用不改变正确与错误轨迹的排序。这只证明给定回报的界，不保证优化必然学会正确策略。

用论文多选配置$\gamma=0.5$、$\rho=0.05$、$C=0.7$、$n_{\max}=6$继续复算。上述$b=0.40$的调用返还$0.5(0.40-0.05)=0.175$；若随后还有一个未通过调用，轨迹价格为$0.175-0.05=0.125$，答对时总回报为1.125。六个调用全部未通过但最终答对，回报仍为0.70；零调用但答错则为0。这里费用单位是训练回报，不是货币或实测运行秒数，部署成本还需另外记录。

策略更新使用组相对策略优化（group relative policy optimization, GRPO），即对同一个问题抽样一组回答轨迹，用组内相对回报决定哪些生成动作应提高概率。CounterCredit把正确性和价格分开处理。若第$j$条轨迹的正确性为$O_j$、价格为$P_j$，组均值为$\overline O,\overline P$，$\sigma_{\min}>0$为避免零分母的标准差下限，论文取0.15，其优势为
\begin{equation}
  \widehat A_j=
    \frac{O_j-\overline O}
         {\max\{\operatorname{std}(O),\sigma_{\min}\}}
       +(P_j-\overline P).
  \label{eq:vision-countercredit-advantage}
\end{equation}
正确性项除以组内标准差，价格项只减均值、不除标准差，所以多一次未核验调用造成的$\rho$差值仍保持原回报尺度。组内全错时，论文把价格优势设为0，避免只根据错误轨迹中的调用数学习。根据优势提高或降低动作概率的策略学习原理可参考\BookChapterRef{chap:rl-learning}{卷七“学习”章}；这里新增的是价格如何进入已有策略更新，而不是另一种视觉编码器。

为了把这个优势接回参数，令$r_{j,k}=\pi_{\bm\theta}(a_{j,k}\mid h_{j,k})/\pi_{\mathrm{old}}(a_{j,k}\mid h_{j,k})$为轨迹$j$在生成位置$k$的概率比。对位置取平均的目标项为$\min\{r_{j,k}\widehat A_j,\operatorname{clip}(r_{j,k},1-\delta,1+\delta)\widehat A_j\}$，$\delta$限制概率比的有效更新幅度；优化器沿增大该项的方向调整策略权重。参考答案、随机图块及探测得分不通过这个目标反向传播，它们先形成停止梯度的优势，真正的梯度来自已抽样动作的对数概率。新的策略再生成下一批轨迹，如此重复到预定训练日程，而不是一直调用到探测得分不再变化。

原实验从Qwen2.5-VL-7B-Instruct的裁剪轨迹监督微调开始，固定视觉编码器与投影器，随后在同一格式和调用预算下做三轮GRPO训练。最终保留学到的策略及裁剪接口，推理时不提供参考答案、不运行这些对照探测；只根据当前状态选择立即回答或实际观察。训练探测每个合格调用还需$K+2$次得分前向，不能把它计作推理免费取得的真值。本文接口可以由后续问题驱动的证据获取回用，但论文当前信号不直接覆盖搜索、OCR文字返回或任意外部工具；随机参照偏差、策略得分误校准和跨模态分布变化也会限制核验的解释。

\subsection{媒体内容的生成}
\label{sec:vision-models-media-generation}

生成一幅雨天街景时，仅有“雨天”这个语义判断还不够，模型需要决定每处颜色、纹理与空间关系。直接在像素中生成可以保留细粒度自由度，但计算量大；在压缩表示中生成可以减少计算，却必须有相应解码器。条件信息怎样进入生成过程，又形成另一项独立选择。

\subsubsection{LDM的空间潜表示与文字条件接入}
\label{sec:vision-models-ldm}

潜空间扩散模型（latent diffusion model, LDM）先学习图像压缩与重建，再在压缩空间中学习生成。这里的潜表示保留二维网格结构，与CLIP把整图汇总成一个语义向量的空间不同。原论文的文字条件LDM通过交叉注意力把文字表示送入去噪网络\scite{pgm-generative-rombach2022}

图像自编码器由编码器$E$与解码器$D$组成。编码得到$\bm Z=E(\bm X)\in\mathbb R^{h'\times w'\times d_z}$，解码得到$D(\bm Z)$。自编码器（autoencoder）通过还原输入来学习压缩坐标，不等于按类别汇总一幅图片。自编码阶段使用重建相关的感知与对抗目标，并配合潜空间正则化，学习在压缩后保留可重建的信息。该阶段完成后固定自编码器，在其潜空间中训练扩散网络，不把图像重建器和每个生成任务从头共同重训。

重建训练的输入与目标都是同一张图：输入先经$E$压缩，再经$D$还原，像素或已有感知网络特征的差异衡量保留信息的误差。感知误差比较的是两图特征，不要求逐像素相同；局部判别器则根据真实与重建图块学习区分，编码器和解码器反过来提高重建图的真实程度。KL正则版本限制潜分布偏离标准高斯的程度，向量量化版本则学习有限代码及其对应，这两种约束不能当作同一个目标。原论文给出不同自编码器设置，选择哪一套便同时确定了后续潜变量尺度和解码器。这个阶段提高重建保真度，不直接训练“根据雨天文字创造新场景”。

为说明计算对象，把潜数组按固定次序展开为向量$\bm z$。取噪声步$k$，给它加入标准高斯噪声$\bm\epsilon$：
\[
  \bm z^{(k)}
  =\sqrt{\bar\alpha_k}\bm z+
    \sqrt{1-\bar\alpha_k}\bm\epsilon\text{，}
\]
其中$\bar\alpha_k$由噪声日程指定，$(k)$表示加噪或去噪过程的步，不是视频帧。文字请求$c$经条件编码器得到$\bm C=E_{\bm\phi}(c)$。去噪网络以当前潜表示、噪声步和文字为条件预测所加噪声，目标为
\begin{equation}
  L_{\mathrm{LDM}}
  =\mathbb E_{\bm X,c,k,\bm\epsilon}
    \left[
      \left\lVert\bm\epsilon-
      \epsilon_{\bm\theta}
      \bigl(\bm z^{(k)},k,E_{\bm\phi}(c)\bigr)
      \right\rVert_2^2
    \right]\text{。}
  \label{eq:vision-ldm-objective}
\end{equation}
原论文的条件设置共同优化$\bm\theta$与$\bm\phi$；后来使用冻结文字编码器的实现，需要另行说明。

这里$\bar\alpha_k=\prod_{j=1}^{k}\alpha_j$，$\bar\alpha_0=1$，$0<\alpha_j<1$，每个$\alpha_j$控制一次保留原状态的比例。独立高斯加噪的合成仍为高斯，所以可以从干净潜表示一次构造任意$k$的带噪输入，而不必真的先算完前$k$步。每次训练先抽图文对，再抽$k$和同维标准噪声，编码干净潜网格后直接计算$\bm z^{(k)}$。网络输出与噪声同形，平方误差对输出的导数为$2(\widehat{\bm\epsilon}-\bm\epsilon)$，优化器把这个差传到去噪网络、文字编码器和连接投影。固定的图像自编码器提供目标坐标，不参与这次更新。以有限批次均值估计式\eqref{eq:vision-ldm-objective}，按所选训练日程和验证结果保留检查点，而不是把抽样期望当成已精确求解的总体目标。

交叉注意力中，查询来自当前去噪网络的空间特征，键和值来自文字序列。假设当前有$n_z$个空间单元、文字有$m$个词元，注意力权重是$n_z\times m$的矩阵，每个空间位置可以按当前状态读取不同文字信息。它没有直接为“红色”指定一个真实对象框，词语与空间的正确对应仍需通过数据和目标学习。

使用时如何从预测噪声得到更干净状态，还需要一步。把带噪关系中的噪声替换成预测值后，先估计干净潜向量：
\begin{equation}
  \widehat{\bm z}
   =\frac{\bm z^{(k)}
       -\sqrt{1-\bar\alpha_k}\,
           \bm\epsilon_{\bm\theta}(\bm z^{(k)},k,\bm C)}
          {\sqrt{\bar\alpha_k}}.
  \label{eq:vision-ldm-clean-estimate}
\end{equation}
这来自对前向加噪等式移项，并非使用时已知真实$\bm z$。若采用标准高斯反向采样，把$\beta_k=1-\alpha_k$和
$\widetilde\beta_k=\beta_k(1-\bar\alpha_{k-1})/(1-\bar\alpha_k)$记为噪声方差系数，反向均值与更新可写为
\begin{equation}
  \begin{aligned}
    \bm\mu_k
      &=\frac1{\sqrt{\alpha_k}}
        \left(\bm z^{(k)}
          -\frac{\beta_k}{\sqrt{1-\bar\alpha_k}}\,
             \bm\epsilon_{\bm\theta}(\bm z^{(k)},k,\bm C)\right),\\
    \bm z^{(k-1)}
      &=\bm\mu_k+\sqrt{\widetilde\beta_k}\,\bm\eta_k,
         \qquad \bm\eta_k\sim\mathcal N(\bm0,\bm I).
  \end{aligned}
  \label{eq:vision-ldm-reverse-interface}
\end{equation}
均值形式由前向高斯链的后验及噪声参数化得到；最后一步不再加随机噪声。不同采样器可以改用确定性或跳步更新，所以这个接口说明基本依赖，不规定所有LDM检查点的实际步数。令$K$为所选噪声日程的末步，且$\bar\alpha_K$足够小，生成便从近似标准高斯的$\bm z^{(K)}$开始，按$k=K,\ldots,1$反复预测，再把末态恢复为二维潜网格，用同一$D$解码，反转训练时的像素尺度处理才得到显示图片。

沿街景例子，假设只查看二维潜网格中的一个两维单元，干净值为$(2,-1)^\top$，抽样噪声为$(0,1)^\top$，$\bar\alpha_k=0.64$。训练输入为$(1.6,-0.2)^\top$；若网络预测噪声$(0.1,0.8)^\top$，误差平方和为$0.01+0.04=0.05$，输出导数为$(0.2,-0.4)^\top$。式\eqref{eq:vision-ldm-clean-estimate}恢复的干净估计为$(1.925,-0.85)^\top$，而不是原来的$(2,-1)^\top$。这条计算明确显示噪声误差怎样改变潜表示；它还不是完整街景像素，因为其他网格位置和解码器共同决定最终图像。

使用时先编码“雨天街道，右侧停着红色汽车”，再从潜空间噪声出发多步去噪，最终用$D$恢复图片。若高宽都压缩为原来的八分之一，潜网格位置数变为$1/64$；这并不意味着全部计算或内存恰好减少64倍，因为通道数、网络宽度和注意力实现仍然不同。压缩还可能损失小字等细节，生成后的事实与控制要求由\ref{chap:vision-image-generation}章的任务协议检查。

\subsubsection{unCLIP的图像表示先验与图像解码}
\label{sec:vision-models-unclip}

LDM直接用文字条件预测空间潜表示，unCLIP则在文字与图像之间加入一个整图语义表示。冻结CLIP后，先验模型根据文字预测CLIP图像表示，另一个生成解码器再依据该表示形成图片。这是DALL-E 2报告中的unCLIP路线，生成器并非CLIP本身\scite{vision-unclip2022}

设文字为$c$，CLIP图像表示为$\bm z_{\mathrm i}$，图像为$\bm X$，模型可写成
\begin{equation}
  p(\bm X\mid c)
  =\int p_{\mathrm{dec}}(\bm X\mid\bm z_{\mathrm i},c)\,
          p_{\mathrm{prior}}(\bm z_{\mathrm i}\mid c)\,
          \mathrm d\bm z_{\mathrm i}\text{。}
  \label{eq:vision-unclip-factorization}
\end{equation}
这里表达模型的分层生成分布，不宣称训练后已等于真实数据分布。先验处理“这段文字可能对应哪些图像语义表示”，解码器处理“这个语义表示可能对应哪些具体像素”。后者还可以读取文字，不能把完整条件一律删成只有图像向量。

训练先验时，图文对中的图像经冻结CLIP形成目标向量。论文比较了两条路线：自回归先验先对图像向量作主成分降维与分量量化，再逐项预测离散编码；扩散先验直接对连续图像向量加噪，并在文字条件下回归干净向量。扩散先验采用干净表示预测目标，不能直接照搬LDM的噪声预测符号而不作转换。

自回归先验先在训练图像表示上估计主成分方向，按方差从大到小保留分量，再把每个分量的值分到离散桶。原报告用319个主成分、每分量1024桶。训练以文字、CLIP文字表示及图文点积的量化条件为前缀，逐分量预测正确桶编号，负对数概率之和训练先验Transformer。使用时先选点积条件，再逐项抽桶、取对应数值、逆主成分变换，恢复候选图像表示。量化与截断产生额外误差，不能把所得向量当作原CLIP图像的精确逆。

扩散先验保留连续向量。令$\bm z_{\mathrm i}^{(k)}$为按前节高斯规则加噪的图像表示，文字及CLIP文字表示合写为条件$c$，先验网络$f_{\bm\psi}$的目标为
\begin{equation}
  L_{\mathrm{prior}}
   =\mathbb E_{\bm X,c,k,\bm\epsilon}
      \left[\left\lVert
         f_{\bm\psi}(\bm z_{\mathrm i}^{(k)},k,c)
           -\bm z_{\mathrm i}\right\rVert_2^2\right].
  \label{eq:vision-unclip-prior-objective}
\end{equation}
目标$\bm z_{\mathrm i}$由冻结CLIP计算；先验预测干净向量，误差传回$\bm\psi$，不改变CLIP坐标。采样时，若需要噪声形式，可用
$\widehat{\bm\epsilon}=(\bm z_{\mathrm i}^{(k)}-\sqrt{\bar\alpha_k}\widehat{\bm z}_{\mathrm i})/\sqrt{1-\bar\alpha_k}$换算，再实施反向更新。论文在使用时采两个候选图像表示，按与CLIP文字表示的点积选择其中一个；这是表示阶段的选择，尚没有生成两幅完整图片。

训练生成解码器时，以真实图像的CLIP表示和相应文字为条件，对带噪图像学习去噪。图像表示被投影到时间步嵌入及额外上下文单元，连接到由GLIDE结构改造的图像生成网络；随后使用图像上采样模块恢复更高分辨率。CLIP在先验与解码器训练期间保持冻结，两个生成模块学习不同的映射。

解码训练的目标现在是像素而非整图向量：图文对给出$\bm X$，CLIP给出$\bm z_{\mathrm i}$，对$\bm X$加噪后用$(\bm z_{\mathrm i},c)$条件预测噪声，学习反向图像分布。先验和解码器分别抽样、计算各自误差并按各自日程更新，最后保留两套权重及上采样模型。生成器先形成$64\times64$图片，再经$256\times256$、$1024\times1024$上采样路径恢复高分辨率；后两级收到较低分辨率图像条件，而不是重新从一段文字独立生成。原文的模块形式改自GLIDE，不表示必须把一个现成GLIDE模型不经训练直接接上即可工作。

对“右侧红车、左侧红色自行车的雨天街景”，先验产生的$\widehat{\bm z}_{\mathrm i}$只表达一组整体条件。固定它和文字、换不同图像噪声，解码器可以产生多幅雨天街景。若两候选表示与文字点积为0.72和0.68，选择第一项只说明这一表示比较规则，不能提前判定最终图片一定有两辆车或位置正确。参考图直接提供图像表示时跳过先验，但仍需像素生成与上采样；要保留原图细节，报告的反演路线还记录扩散末端状态等残余信息，而不是仅保存一个CLIP向量。

这一分层路径把语义条件与像素细节分开：相近图像表示可以对应纹理、布局细节不同的结果，表示匹配并不决定完整重建。参考图片的变化可以展示哪些内容在生成中较易保持，但“语义相似”仍不能替代对原图每处细节的检查。

\subsubsection{GLIDE的文字条件与语义引导}
\label{sec:vision-models-glide}

条件输入规定去噪网络能够读取什么，引导规定使用时怎样组合或修正预测，生成后重排则只决定从已有候选中选谁。GLIDE在像素空间学习文字条件扩散，提供了比较这三个环节的具体入口。本节建立生成接口，后续图像创作在这个接口上选择控制方式与评价协议\scite{vision-glide2022}。

GLIDE先用文字Transformer编码请求，将文字信息接入图像去噪网络。训练根据图文配对构造带噪图像，让去噪预测依赖文字。为了实现无分类器引导，训练时还以一定概率去掉文字条件，使同一个网络能给出条件预测和无条件预测。使用时对同一噪声状态分别计算两者，再按选定强度组合。

一条训练记录仍是图像与描述，但带噪对象改为像素数组。文字网络的末端表示接入噪声步条件，文字序列又作为去噪网络注意力层的键和值，当前图像空间特征提供查询。抽样噪声步后，以所加噪声为目标，用平方误差训练文字编码、去噪网络及连接权重，论文还训练反向方差相关项。文字条件随机置为空序列时，网络不能再读取原描述；所以无条件预测是同一网络在真实训练过的空条件下运行，并非把预测中的某个词向量事后删掉。

令$\bm\epsilon_{\mathrm c}$与$\bm\epsilon_{\varnothing}$分别表示同一带噪图片、同一噪声步上的条件与空条件预测，无分类器引导（classifier-free guidance, CFG）用
\begin{equation}
  \widehat{\bm\epsilon}
     =\bm\epsilon_{\varnothing}
        +g(\bm\epsilon_{\mathrm c}-\bm\epsilon_{\varnothing})
  \label{eq:vision-glide-cfg-interface}
\end{equation}
替代反向均值计算中的噪声预测，$g$是使用时设定的无量纲强度。$g=1$恢复条件预测，$g>1$沿条件与空条件预测的差再向外推；这里的线性组合作用于噪声预测，不是把两幅完成图片混合。通过式\eqref{eq:vision-ldm-reverse-interface}同类的像素反向更新形成下一状态，重复到无噪声末态，再由上采样网络扩大分辨率。训练的步数和优化停止由原模型日程指定，推理的采样步数则由所选采样器指定，二者不是同一个计数。

固定街景文字和一个二维教学噪声状态，假设$\bm\epsilon_{\varnothing}=(0.2,0.4)^\top$、$\bm\epsilon_{\mathrm c}=(0.1,0.8)^\top$。取$g=2$得到$(0,1.2)^\top$。这个结果不是条件预测的概率，而是要送入下一反向更新的数值；第二维已经超过两种原预测的范围，说明引导是外推。增加强度可能改变条件遵循与样本多样性，也可能引入异常，不能由组合形式推出越大越好。

噪声感知CLIP引导需要另一套图文比较模型。它必须读取当前噪声水平的图像，把文字相似度对图像的梯度用于修正生成方向。原GLIDE实验专门训练了适应噪声输入的CLIP，不能把只看清晰图片的原CLIP不加说明地放在同一位置。固定请求与初始噪声时，改变引导强度会改变每一步的预测或更新，最终像素也可能随之改变。

设噪声适应图像编码与文字编码形成匹配分数$s_{\bm\phi}(\bm x^{(k)},c)$，反向图像分布的原均值为$\bm\mu_k$、协方差为$\bm\Sigma_k$，CLIP梯度引导写成
\begin{equation}
  \widehat{\bm\mu}_k
   =\bm\mu_k+g_{\mathrm{clip}}\bm\Sigma_k
           \nabla_{\bm x^{(k)}}s_{\bm\phi}(\bm x^{(k)},c).
  \label{eq:vision-glide-clip-interface}
\end{equation}
梯度指出当前像素怎样局部变化会提高匹配得分，协方差把它转成与反向更新相容的均值修正。CLIP权重在采样中固定，求导对象是当前图像状态，不是临时微调编码器。教学上若$\bm\mu_k=(0.5,0.5)^\top$、$\bm\Sigma_k=0.1\bm I$、梯度为$(1,-2)^\top$、$g_{\mathrm{clip}}=2$，修正均值就是$(0.7,0.1)^\top$。它发生在本步抽样之前，因而能改变未来轨迹；生成后的CLIP排序没有这条像素梯度路径。

生成后重排则先完成$K$个候选，之后才计算各自的匹配分数并选择。重排不修改任何候选，也不能从全部错误的候选中创造正确答案。三种环节可以组合，但应把额外模型调用、梯度和候选数量计入预算。图文相似度仍需与对象数量、属性归属和空间关系的独立核验配合。

\subsubsection{流匹配的速度目标与连续生成}
\label{sec:vision-models-flow-matching}

媒体生成不一定把网络输出解释为“应去掉多少噪声”，也可以直接学习“当前表示应向哪个方向变化”。\term{流匹配}（flow matching）就是通过回归预先构造路径上的速度来训练连续生成模型。它把训练时可计算的局部目标，与使用时求解整条生成轨迹联系起来\scite{vision-flowmatching2023}

用一个明确的教学构造说明。设数据潜向量为$\bm z_1\in\mathbb R^d$，噪声为$\bm z_0\sim\mathcal N(\bm0,\bm I)$，路径参数$s\in[0,1]$。在两者之间定义直线路径
\begin{equation}
  \bm z(s)=(1-s)\bm z_0+s\bm z_1,
  \qquad
  \bm u=\frac{\mathrm d\bm z(s)}{\mathrm ds}
       =\bm z_1-\bm z_0\text{。}
  \label{eq:vision-flow-path}
\end{equation}
$s=0$对应噪声，$s=1$对应数据，$s$是无量纲的生成进度，速度的单位是每单位生成进度的潜坐标变化，不是米每秒或视频帧速度。这里采用无残余噪声的理想化端点以便手算；原流匹配论文也讨论终点保留微小方差等路径，不能把所有实现的日程等同于这一个式子。

训练时取一对数据与条件$(\bm z_1,c)$，抽样噪声$\bm z_0$和时刻$s$，就能直接算出中间状态$\bm z(s)$与目标速度$\bm u$。速度网络$\bm v_{\bm\theta}$的平方误差目标为
\begin{equation}
  L_{\mathrm{FM}}
  =\mathbb E_{\bm z_1,c,\bm z_0,s}
     \left[
       \left\lVert
         \bm v_{\bm\theta}(\bm z(s),s,c)
           -(\bm z_1-\bm z_0)
       \right\rVert_2^2
     \right]\text{。}
  \label{eq:vision-flow-loss}
\end{equation}
计算一次训练目标不必先求解从0到$s$的微分方程，因为中间状态有显式构造。更新的是速度网络参数，样本端点和抽样噪声给出监督。

有限训练集上，把$n$个样本的随机路径误差取平均，便得到一次批次目标。若某个输出速度的残差为$\bm r=\bm v_{\bm\theta}-\bm u$，平方误差导数为$2\bm r$，再经速度网络反传到可训练权重。可从已有生成检查点初始化，也可随机初始化，具体依实现而定；媒体编码器若固定，目标端点在训练中不因这次速度误差改变。优化器按有限训练日程更新，保留速度网络、条件编码和对应媒体解码器。流匹配提供目标构造，不给任意网络和数据配置一个自动收敛到真实分布的保证。

使用时却不知道目标$\bm z_1$，所以不能直接取$\bm z_1-\bm z_0$走向答案。模型需要从一个新噪声出发，根据当前状态预测速度，求解
\begin{equation}
  \frac{\mathrm d\widehat{\bm z}(s)}{\mathrm ds}
    =\bm v_{\bm\theta}(\widehat{\bm z}(s),s,c),
  \qquad
  \widehat{\bm z}(0)\sim\mathcal N(\bm0,\bm I)\text{，}
  \label{eq:vision-flow-ode}
\end{equation}
再把$\widehat{\bm z}(1)$送入媒体解码器。简单的欧拉法在$s_k$处用
$\widehat{\bm z}^{(k+1)}
=\widehat{\bm z}^{(k)}+(s_{k+1}-s_k)
\bm v_{\bm\theta}(\widehat{\bm z}^{(k)},s_k,c)$
更新。有限步长有数值误差，网络也有预测误差，两者不能混为同一项。

为什么从许多不同样本路径学习，仍能得到一个只依赖当前状态的速度？固定输入$(\bm z,s,c)$时，平方损失的最优预测是目标速度的条件均值。令$\bar{\bm u}=\mathbb E[\bm u\mid\bm z,s,c]$，则
\[
  \mathbb E[\lVert\bm v-\bm u\rVert_2^2\mid\bm z,s,c]
  =\lVert\bm v-\bar{\bm u}\rVert_2^2+
    \mathbb E[\lVert\bm u-\bar{\bm u}\rVert_2^2\mid\bm z,s,c]\text{。}
\]
展开平方后，交叉项因$\mathbb E[\bm u-\bar{\bm u}\mid\bm z,s,c]=\bm0$消失。右侧第二项不依赖$\bm v$，因此网络学习的是穿过当前状态的样本路径速度的条件平均。在适当正则条件下，这一边缘速度场可以推动对应的概率分布路径；它不要求网络辨认训练时某条唯一的端点连线。

这个分布结论还需要速度场及密度足够规则，使样本运动可由连续性方程描述。对条件路径密度$p_s(\bm z\mid c)$，质量守恒写为$\partial_s p_s=-\nabla_{\bm z}\cdot(p_s\bar{\bm u})$；局部流入流出的差决定该位置密度如何改变。理论上的条件平均场满足这条对应关系，而有限数据和有限容量训练得到的$\bm v_{\bm\theta}$只是近似。因此，训练时知道直线端点，推理时只知道当前状态；两者不应被画成同一条已知答案路径。

取二维教学样本$\bm z_0=(0,1)^\top$、$\bm z_1=(2,-1)^\top$，目标速度为$(2,-2)^\top$，$s=0.25$时状态为$(0.5,0.5)^\top$。假如使用时在该位置预测速度$(1.2,-2.8)^\top$，步长0.25给出下一状态$(0.8,-0.2)^\top$，而该训练直线在$s=0.5$处为$(1,0)^\top$。这个差异说明一次速度预测误差怎样传到位置，不能用单条训练直线宣称网络推理必然沿直线运行。图\ref{fig:vision-flow-step}给出几何对应。

\begin{figure}[htbp]
  \centering
  % 直线路径为教学构造；粗箭头使用特意偏离目标的速度预测。
\begin{tikzpicture}[x=3cm,y=1.65cm,font=\figurefont\small,
  arr/.style={-{Stealth[length=2mm]},BookInk,line width=0.8pt}]
  \draw[arr] (-0.15,-1.4) -- (2.55,-1.4) node[right] {$z_1$};
  \draw[arr] (-0.15,-1.4) -- (-0.15,1.5) node[above] {$z_2$};
  \foreach \x in {0,1,2}{
    \draw[BookMuted,line width=0.5pt] (\x,-1.4) -- (\x,-1.46);
    \node[below] at (\x,-1.46) {$\x$};
  }
  \foreach \y in {-1,0,1}{
    \draw[BookMuted,line width=0.5pt] (-0.15,\y) -- (-0.19,\y);
    \node[left] at (-0.19,\y) {$\y$};
  }
  \draw[BookBlue,line width=0.9pt] (0,1) -- (2,-1);
  \fill[BookBlue] (0,1) circle[radius=1.1mm];
  \fill[BookTeal] (2,-1) circle[radius=1.1mm];
  \node[anchor=south west] at (0,1.13) {$s=0,\ (0,1)$};
  \node[anchor=west] at (2.02,-0.96) {$s=1$};
  \fill[BookInk] (0.5,0.5) circle[radius=1mm];
  \node[anchor=south west] at (0.52,0.62) {$s=0.25,\ (0.5,0.5)$};
  \draw[BookTeal,line width=1pt,fill=white] (1,0) circle[radius=1.1mm];
  \node[anchor=west] (ideal) at (1.35,0.1) {路径上$(1,0)$};
  \draw[BookTeal,line width=0.6pt] (1.04,0) -- (1.3,0.1);
  \draw[-{Stealth[length=2.2mm]},BookAmber,line width=1.3pt]
    (0.5,0.5) -- (0.8,-0.2);
  \fill[BookAmber] (0.7825,-0.23) rectangle (0.8175,-0.17);
  \node[anchor=east] at (0.76,-0.35) {预测$(0.8,-0.2)$};
  \draw[BookMuted,dashed,line width=0.8pt] (0.82,-0.17) -- (0.97,-0.03);
  \node[anchor=west] at (0,-2.05) {粗箭头：$0.25(1.2,-2.8)^\top$};
\end{tikzpicture}

  \caption{流匹配教学样本的路径与一次欧拉更新。实线连接已知训练端点，粗箭头为预测速度乘以步长后的位移，虚线标示与训练路径上$s=0.5$位置的差别。横纵轴是两个潜变量坐标；路径参数$s$不是视频中的事件时间。}
  \label{fig:vision-flow-step}
\end{figure}

即便每条条件路径都很简单，整体数据分布对应的轨迹也未必是直线，更不能仅据此声称得到了全局最优传输。规范化流（normalizing flow）描述通过可逆变换建立分布的模型，流匹配描述训练速度场的一类目标，两者相关但不是同义词。若$\bm z_1=T(\bm z_0)$为可微可逆变换，密度还须补上体积变化：
\begin{equation}
  \log p_1(\bm z_1)
    =\log p_0(\bm z_0)
       -\log\left|\det
          \frac{\partial T(\bm z_0)}{\partial\bm z_0}\right|.
  \label{eq:vision-flow-density-transform}
\end{equation}
例如$T(z)=2z$把同样概率质量铺在两倍长度上，所以密度减半，对数密度减去$\log2$。这解释了雅可比行列式为何不能从概率计算中省略。

连续情形把许多微小变换相接。若速度场对空间可微且产生唯一可逆轨迹，沿轨迹的密度变化为
\begin{equation}
  \frac{\mathrm d}{\mathrm ds}
        \log p_s(\widehat{\bm z}(s))
    =-\nabla_{\bm z}\cdot
          \bm v_{\bm\theta}(\widehat{\bm z}(s),s,c).
  \label{eq:vision-flow-density-divergence}
\end{equation}
散度是速度对各同名坐标导数之和，度量局部体积的扩张或收缩；扩张会降低单位体积上的密度。该式来自连续性方程和沿轨迹的链式求导，不是流匹配平方误差中的另一项必选监督。只需生成样本的系统可以不计算似然，要求密度的系统还要积分散度并计入求解代价。相关条件与变换原理可参考\BookChapterRef{chap:pgm-generative}{卷六“生成模型”章}。

\subsubsection{LTX-2的音视频双分支与联合解码}
\label{sec:vision-models-ltx2}

若只生成敲门画面，再由另一个模型补声音，后一步只能适应已固定的动作。联合生成可以让画面和声音在生成过程中互相影响。LTX-2技术报告v1采用分别表示、相互连接的音视频分支，学习文字条件下的联合媒体输出。

视频经因果变分自编码器（variational autoencoder, VAE）编码为空间与时间潜网格，音频经频谱转换和独立音频VAE编码为时间潜序列。VAE把媒体压缩到受分布约束的连续表示，再由对应解码器还原；这里的因果约束限制编码使用未来时间的信息，不等于自由问答的推理因果性。两个分支的长度、通道宽度和压缩方式可以不同，因此不需要把音频潜向量误解释成视频像素。文字经编码与连接模块形成条件，分别进入两种媒体的生成分支。

每个分支具有自己的注意力和前馈计算，又通过双向交叉注意力交换信息。视频查询读取音频的键和值，音频查询反过来读取视频的键和值。视频分支内部同时使用时空位置，跨音视频连接则着重使用时间位置，因为音频时间单元没有图像平面上的横纵坐标。时间编码使对应关系可被表达，但正确的敲击同步仍要从配对数据和生成结果检验。

训练把配对音画编码为各自目标潜表示，构造噪声状态，用流匹配监督两个分支的速度\scite{vision-ltx22026}
沿用前节的教学记号，可以把联合状态写成$(\bm z^{\mathrm v}(s),\bm z^{\mathrm a}(s))$，两边的速度预测都依赖当前音画状态及文字条件。各项损失的尺度需要按维度和具体训练设置处理，不能不计音视频长度差异而把元素误差总和直接比较。

报告说明结构与条件路径，公开训练源码进一步给出噪声与速度的符号方向：它以$\sigma=1$表示噪声、$\sigma=0$表示干净数据，构造$(1-\sigma)\bm z_{\mathrm{clean}}+\sigma\bm\epsilon$，目标为$\bm\epsilon-\bm z_{\mathrm{clean}}$。令本章$s=1-\sigma$，便恢复式\eqref{eq:vision-flow-path}的噪声到数据方向，同时速度需变号；不能在减小$\sigma$的采样器中误用本章正向速度。源码中的音视频批次共用$\sigma$，各自抽同形噪声，各自按有效元素平均平方误差，最后把两项相加；作为条件的已知视频位置可由掩码排除监督。这个可核查实现不代表已公开全部基础预训练日程。

文字条件来自冻结的Gemma3多层隐表示，经特征提取投影及每模态的文字连接器处理。报告的初期阶段共同训练特征投影与生成模型，之后固定该投影；文字连接器仍与音视频生成模块共同训练。两个媒体自编码器和声码器提供已定义的解码空间，生成训练保留它们的输入输出约定。学习后的速度、连接模块、两套媒体解码器及文字处理器一起构成使用接口，不能只保留一个视频分支就宣称仍有联合声音输出。

若把一个敲门教学片段的状态简化为视频和音频各一个标量，噪声端为$(0,1)$、数据端为$(2,-1)$，目标正向速度为$(2,-2)$。$s=0.25$时联合状态为$(0.5,0.5)$；若两分支预测为$(1.5,-1.5)$，总平方残差为$0.25+0.25=0.5$。步长0.25更新到$(0.875,0.125)$。真实系统有长时间潜序列与双向读取，不能从这两个标量直接判断敲门同步；例子只显示两种媒体在同一次生成步中各有目标、误差和更新。

使用时，从两个潜空间的噪声开始共同更新。视频潜表示由视频VAE解码为帧，音频潜表示先恢复梅尔频谱，再由声码器恢复波形。若文字要求“人物敲门两次，随后停下”，验收至少要对应两次接触、两次声音以及停顿，不能只听到类似敲门声就认定通过。联合生成不同于为固定视频配声，接口本身也不提供视觉问答；LTX-2中的文字编码器服务生成条件，不等于系统具有自由问答输出头。

\section{统一理解与生成模型}
\label{sec:vision-models-unified}

前面的方法可以形成可比较表示、语言或媒体。如果用户希望“解释这张图，再按要求画一个变体”，可以在一个主干中共享计算，也可以把理解模型和生成器连接起来。判断“统一”的含义，需要看哪些参数、表示和训练目标实际共用，不能只看界面是否只有一个输入框。

\subsection{Show-o2的共享视觉表示与双输出头}
\label{sec:vision-models-showo2}

理解图像需要语义特征，生成图像还需要纹理和局部结构。Show-o2在同一视觉潜空间中提取这两类信息，送入共享语言主干，再用不同输出头分别处理文字和视觉生成。本节依据论文v3的结构及训练说明\scite{vision-showo22025}

图像和视频由3D因果VAE编码。因果性约束编码中未来信息的使用，图像可以作为相应的单帧情形处理。视觉潜表示经过两条路径：语义层从潜网格中提取高层信息，投影层保留便于后续处理的低层特征。语义层在主要两阶段训练之前，利用预训练SigLIP的图像特征进行蒸馏，使其能够从干净和带噪潜表示中提取语义。

两个路径按对应的空间及时间位置对齐，沿特征维拼接，再经归一化与MLP融合为共享视觉表示。这里的共享表示保留多个位置，并非把整幅图压成一个CLIP向量。干净参考图以路径终点$s=1$的状态进入，生成中的视觉块则带有当前路径时刻。相同的视觉接入模块因此需要处理不同噪声水平。

设语义路径输出为$\bm S\in\mathbb R^{n_z\times d_{\mathrm s}}$、低层路径输出为$\bm P\in\mathbb R^{n_z\times d_{\mathrm p}}$，融合输入为$[\bm S,\bm P]\in\mathbb R^{n_z\times(d_{\mathrm s}+d_{\mathrm p})}$。同一行须来自相同潜网格范围，再经均方根归一化和两层映射得到语言宽度的$\bm U$。归一化调整每行的数值尺度，拼接则保留两类特征供映射共同读取；行数不匹配时不能只因同属一张图就直接拼接。3D VAE在这里把时空像素压缩为可连续生成的潜网格，实际使用的权重来自Wan2.1，语言主干来自Qwen2.5-Instruct。语义路径由预训练SigLIP视觉层与新的潜块嵌入组成，先接受干净图片特征的蒸馏，再接入主要训练。

文字嵌入和视觉块按交错顺序进入主干。序列块之间遵守因果顺序，使后续输出可以依赖先前条件；同一视觉块内部允许完整注意力，使一个待生成图像或视频块中的位置共同参与当前速度预测。这里的“因果”描述输出块的条件顺序，不等于每个视频帧只能读取此前帧。

语言头把隐状态转换为下一词元分布，流头则预测视觉潜表示的速度。后者包含带生成时刻调制的Transformer层，其输出经过数值更新和VAE解码才成为媒体。两种训练目标可以概括为
\begin{equation}
  L=\lambda_{\mathrm{text}}L_{\mathrm{NTP}}+L_{\mathrm{FM}}\text{，}
  \label{eq:vision-showo-loss}
\end{equation}
其中$L_{\mathrm{NTP}}$是文字的下一词元预测损失，$L_{\mathrm{FM}}$是视觉速度误差，$\lambda_{\mathrm{text}}$控制文字项权重。两项可以通过共享主干联系起来，但输出维数和误差对象不同。

一条理解样本把干净图像块和问题作为条件，监督后续答案文字；一条生成样本把文字描述作为条件，对目标图像潜表示抽噪声和路径时刻，监督流头速度。交错样本还可以让先前干净图像与文字成为后来图像的条件，不应对作为条件的图片强行添加“预测下一词元”的目标。训练以各样本允许的目标位置分别计算式\eqref{eq:vision-vlm-text-objective}和式\eqref{eq:vision-flow-loss}，再按式\eqref{eq:vision-showo-loss}合并。目标权重改变共享参数受到两种梯度的比例，不能通过一个标量权重把词元误差和潜变量误差解释成同一种单位。

流头每次给出与视觉潜网格相同形状的速度，但语言头仍给出词表维数的logit。比如一个$27\times27$的视觉网格有729个空间单元，生成头对各潜块预测若干连续分量；语言头在一个回答位置对每个词元编号计分。这一差别解释了“统一”为什么不要求相同输出头。对共享权重$\bm\theta_{\mathrm{shared}}$，一次联合更新的梯度是$\lambda_{\mathrm{text}}\nabla L_{\mathrm{NTP}}+\nabla L_{\mathrm{FM}}$；两个梯度可能指向不同方向，因此共享可以传递信息，也可能产生任务冲突，不能仅据共享结构保证两项能力同时改善。

原论文的第一阶段只训练投影、时空融合和流头，保留已有语言知识，并学习视觉生成接口；第二阶段用高质量理解与生成数据适配除VAE外的整体模型。图\ref{fig:vision-showo-heads}把共享范围与两个输出头分开表示。VAE不随这两个阶段更新，统一训练也不意味着每个部件都在每个阶段改变。

具体更新从前述检查点与蒸馏后的语义层开始。第一阶段对图文批次重算文字预测及潜速度误差，只改变新接入的投影、融合与流头；冻结语言层仍需要把误差传到连接输入。第二阶段再放开除VAE外的模型，用高质量理解和生成样本共同更新，并按有限步数训练日程保留最终检查点。原论文1.5B设置以AdamW训练，并在第一阶段和第二阶段改变文字项权重；视频与交错版本另外加入相应样本。已有权重来源、目标位置、开放参数及数据混合共同决定实际得到的模型，不能只说“统一训练”而省略这四项。

\begin{figure*}[htbp]
  \centering
  \begin{tikzpicture}[x=1cm,y=1cm,font=\figurefont\small,
  block/.style={draw=FigureModel,fill=FigureModel!10,line width=0.8pt,
    minimum height=9mm,minimum width=20mm,align=center},
  input/.style={block,draw=FigureInput,fill=FigureInput!10},
  output/.style={block,draw=FigureOutput,fill=FigureOutput!10},
  arr/.style={-{Stealth[length=2mm]},BookInk,line width=0.8pt,
    rounded corners=1.2mm}]
  \node[input] (z) at (1.1,1.8) {干净／带噪\\视觉潜表示};
  \node[block] (sem) at (4.0,2.7) {语义层};
  \node[block] (low) at (4.0,0.9) {低层投影};
  \node[block] (fuse) at (6.9,1.8) {按位置拼接\\与融合};
  \node[block] (trunk) at (9.8,1.8) {共享主干};
  \node[block] (text) at (12.7,2.7) {语言头};
  \node[block] (flow) at (12.7,0.9) {流头};
  \node[output] (words) at (15.4,2.7) {回答词元};
  \node[output] (media) at (15.4,0.9) {更新潜状态\\VAE解码};
  \draw[BookInk,line width=0.8pt] (z.east) -- (2.55,1.8);
  \draw[arr] (2.55,1.8) |- (sem.west);
  \draw[arr] (2.55,1.8) |- (low.west);
  \draw[arr] (sem.east) -- (5.45,2.7) |- (fuse.west);
  \draw[arr] (low.east) -- (5.45,0.9) |- (fuse.west);
  \draw[arr] (fuse) -- (trunk);
  \draw[BookInk,line width=0.8pt] (trunk.east) -- (11.25,1.8);
  \draw[arr] (11.25,1.8) |- (text.west);
  \draw[arr] (11.25,1.8) |- (flow.west);
  \draw[arr] (text) -- (words);
  \draw[arr] (flow) -- (media);
  \node (prompt) at (9.8,0.05) {文字上下文};
  \draw[arr] (prompt) -- (trunk);
  \node[anchor=west] at (0,-0.75) {视觉块之间保持因果顺序，块内位置允许完整注意力。};
\end{tikzpicture}

  \caption{Show-o2的共享视觉接入与双输出头示意，依据论文v3重绘。视觉潜表示的语义路径和低层路径按位置融合，再与文字一起进入共享主干。语言头给出词元，流头给出视觉潜变量的速度；生成分支需迭代更新潜状态并重新经过模型，图中省略迭代回边以突出共享和专用模块的边界。}
  \label{fig:vision-showo-heads}
\end{figure*}
\FloatBarrier

对于“解释参考图，再生成雨天变体”，参考图先经过干净视觉路径，语言头生成解释；随后的视觉块依据此前上下文，通过流头逐步形成媒体。解释正确并不自动保证生成正确，生成中仍可能丢失刚才已识别的对象关系。论文包含文字、图像和视频路线，但视频与交错数据的训练安排随变体而异，报告中的较大7B变体没有沿用同样的视频生成训练数据。选择检查点时必须确认实际训练范围，不能把模型族的全部展示合并成每个版本的能力，更不能额外推断音频生成。

两条输出路径可以继续在街景上复算。理解时，参考图按$s=1$送入，文字头逐词给出“右侧是红色汽车”；生成开始标记出现后，系统分配新的视觉噪声块，保留此前图片和文字作为条件。假设新块的一个二维单元在$s=0.25$为$(0.5,0.5)^\top$，流头预测$(1.2,-2.8)^\top$，步长0.25更新为$(0.8,-0.2)^\top$。更新后的块必须重新经过双路径融合、共享主干与流头，才得到下一时刻速度，不能把第一次隐表示重复使用到终点。完整视觉块到$s=1$后经VAE解码成图片，再作为后续交错输出的上下文；词元生成和潜空间求解由不同停止条件控制。

\subsection{理解模型与专用生成器的接口组合}
\label{sec:vision-models-composition}

另一条路线是由视觉语言模型读取图片和请求，形成明确的内容说明，再调用专用生成器。以“把玩具狗旁的花瓶改成蓝色，其余尽量保持”为教学请求，理解模型可以输出目标短语“花瓶”、保留对象“玩具狗”和修改属性“蓝色”。若生成器支持局部编辑，还需要把目标短语转换成原图坐标中的区域或掩码，并把原图作为独立条件传递。

只传一句“蓝色花瓶旁有玩具狗”时，生成器没有收到原图的精确纹理和位置，因而只能依据语义重建场景。传原图与掩码可以保留更多空间条件，但若理解模型把目标区域选成玩具狗，生成器仍可能准确执行了错误区域上的请求。组合系统应保存中间说明、区域来源、输入图片和最终结果，使错误能定位到读取、传递还是生成环节。

连续特征也可以通过适配模块传给生成器，减少必须把全部细节说成文字的限制。然而，不同编码器的特征坐标和尺度未必兼容，适配模块通常需要自己的训练依据。分别训练后直接连接、固定两端学习适配器、共同更新多个模块，是不同的训练安排，不能仅因推理时连续调用就称为共同学习。

与Show-o2相比，这类组合通常不共享理解与生成主干，可以独立替换生成器或读取模型，却增加了中间格式转换和多次调用。共享主干减少某些重复计算，但也需要在同一组参数中协调任务。两种路线的选择取决于能传递哪些证据、哪一步容易出错，以及最终成果是否满足请求。

\subsection{理解、图像生成与反思编辑的共同训练}
\label{sec:vision-models-unireason}

共享理解与生成接口以后，还可能遇到两个不同缺口。用户要求“从对面看同一条街，两辆车的左右关系随视角改变”，字面条件需要转成明确视角；初始图片即使依据正确文字生成，仍可能把两车画在错误位置。前一个问题发生在生成前的条件形成，后一个问题需要看见已生成图片之后才能判断。2026年预印本UniReason 1.0把生成前的知识推理和生成后的反思编辑放入共同训练，所核查版本为v1，不把其有限一轮修订推广成任意次数自我改进\scite{vision-unireason-2026}。

UniReason基于Bagel已有检查点，而不是将Show-o2直接追加一个反思头。Bagel的Transformer混合结构（mixture of transformers, MoT）使理解与生成专家在共同的交错上下文中计算：理解路径读取视觉编码特征并预测文字，生成路径在VAE潜空间预测连续速度。这里“专家”指面向不同输入与目标的参数路径，不是外部人工审稿者。模型有统一上下文和可交互表示，但并非理解与生成全部权重相同；语言预测和图像速度也仍有不同输出目标。已有ViT提供观测特征，VAE提供生成坐标，本节集中解释它们怎样在交错样本中接受监督，不另把Bagel展开成一套任务算法。

生成前的知识推理样本应包含请求$c$、目标文字推理$T$与目标图像$\bm I$；编辑样本还包含原图$\bm I_{\mathrm{src}}$。例如请求涉及相反视角时，$T$把场景对象、观察方向和相应左右关系写成可用的生成条件，而不是仅重复“画得正确”。原论文的文生图数据从人工选择的知识种子和类别说明开始，由Gemini-2.5 Pro扩充请求并生成文字推理，再由Qwen-Image形成图片；编辑数据以UniREdit-Data-100K的原图、指令、目标结果为基础补充推理。Gemini另按指令符合、视觉质量和推理正确性筛选样本。人工种子、已有记录、合成教师和筛选器承担不同职责，筛选通过仍不等于每一条合成知识都已成为事实真值。

反思编辑样本又有不同的监督边界。初始生成器先给出$T^{(0)}$与草图$\bm I^{(0)}$；视觉核验器比较图片与请求，形成可执行的偏差说明$T^{(1)}$；编辑教师根据原草图和偏差说明形成$\bm I^{(1)}$。论文用Gemini-2.5 Pro核验对象是否存在、属性是否正确以及风格、真实程度和审美方面的偏差，由Qwen-Image-Edit实施编辑，再用比较判断保留被认为改善且符合反馈的结果。一条记录因此保存请求、初始文字、初始图片、反思文字及修订图片，不能只存最终图片后声称模型学过“观察错误”的过程。

训练分两阶段。第一阶段固定基础模型的理解分支，只更新生成分支，用没有推理文字的已有文生图和编辑数据加强图像条件落实；它不监督模型生成解释。第二阶段从这一检查点继续，开放模型的理解与生成参数，用知识推理和反思编辑样本共同微调。单轮知识样本同时监督$T$和$\bm I$；反思样本把$T^{(0)},\bm I^{(0)}$作为上下文，不对这段草图过程计算目标，而只监督$T^{(1)},\bm I^{(1)}$。这避免把待纠正草图也当成应完整模仿的输出。开放参数只说明哪些权重允许变化，真正的梯度仍由受监督的位置及可见条件决定。

把文字目标位置记为$\mathcal T$、被监督图像位置记为$\mathcal G$，第$g$幅目标图片的潜向量为$\bm z_{g,1}$，抽样噪声为$\bm z_{g,0}$，其可见上下文为$C_g$。按本章噪声到数据的方向记号，两个误差可以写成
\begin{equation}
  \begin{aligned}
    L_{\mathrm{text}}
      &=-\sum_{j\in\mathcal T}
           \log p_{\bm\theta}(y_j\mid C_j),\\
    L_{\mathrm{img}}
      &=\sum_{g\in\mathcal G}\mathbb E_{s,\bm z_{g,0}}
        \left\lVert
          \bm v_{\bm\theta}(\bm z_g(s),s,C_g)
              -(\bm z_{g,1}-\bm z_{g,0})
        \right\rVert_2^2,\\
    L_{\mathrm{joint}}
      &=\lambda_{\mathrm{text}}L_{\mathrm{text}}
         +\lambda_{\mathrm{img}}L_{\mathrm{img}}.
  \end{aligned}
  \label{eq:vision-unireason-joint-objective}
\end{equation}
这里$C_j$只含当前目标位置之前允许读取的文字与图片，图像项则使用式\eqref{eq:vision-flow-path}的可构造中间状态。若实现用反向噪声参数，需如前节同时变换时间和速度符号。上式是按本章记号整理的监督关系，不声称论文给每个像素单独定义了知识正确性损失。文字负对数概率训练参考反思措辞，图像速度误差训练目标媒体，数据筛选才提供两者是否符合知识及请求的外部依据。

实际批次把不同类型记录打包，按监督掩码累积两种误差，再反传到各自专家和参与条件计算的模块。原论文第二阶段取$\lambda_{\mathrm{text}}=2$、$\lambda_{\mathrm{img}}=1$，用Adam和学习率预热、余弦日程；第一阶段30K次迭代，第二阶段10K次迭代。这里开放的是统一模型的理解与生成分支，公开训练说明仍将已有VAE设为冻结，不能把“全模型微调”解释成重训媒体坐标。迭代次数规定何时停止该实验，不保证推理已经正确或每幅图都会越改越好。最终保留经过共同微调的统一模型及对应视觉处理、VAE解码接口；合成教师和筛选器是训练数据构造者，不是每次推理必须另调用的隐藏部件。

沿街景例子作一次可检查的教学推演。先固定几何简化：只有观察方向改变，画面横向坐标映射为$x'=w-x$，不涉及遮挡及透视。原图汽车中心横坐标为900、自行车为300，宽度$w=1200$，新视角条件应是汽车在300、自行车在900。模型先生成这段明确条件$T^{(0)}$，再从噪声经流更新得到$\bm I^{(0)}$。假设草图仍把汽车放在右侧，模型重新编码草图，反思文字指出“汽车与自行车的左右位置仍按原视角，应交换”，然后把草图和这条修订条件送入编辑路径。公式中的文字目标是这条反思，图像目标是符合新左右关系的$\bm I^{(1)}$；初始错误草图只供读取，不再被标成应模仿的正确图像。

再看这个样本的一次目标计算。若反思中的某目标词元概率为0.5，其负对数概率为0.6931；假设被监督修订图像的一个二维潜单元残差为$(0.1,-0.2)^\top$，平方误差为0.05。仅计这两个教学单元，权重组合为$2\times0.6931+0.05\approx1.4363$。实际目标还要累积其他词元、潜单元及批次，不可把这个数冒充论文训练日志。它说明文字与图像误差怎样在同一次参数更新中共同起作用，也说明“共同训练”不是仅在推理时先后调用两个模型。

推理没有教师目标图像。模型先根据请求及可选原图形成文字条件，再生成草图；若启用反思路径，就以草图为新的可读视觉内容形成偏差说明，并执行一轮编辑，保留初始与修订结果供比较。修订图像不是将原噪声继续多走一步，而是带有草图条件的另一次图像生成过程，需额外编码、文字预测、流求解和解码。论文实现的一轮细化既可能纠正位置，也可能损坏原来正确的颜色或对象；验收应同时检查原条件保持与新增修改。真正相反视角的场景还会产生遮挡与透视变化，前面的$x'=w-x$只用于教学复算，不能当作一般三维视角转换公式。后续图像创作章可回用本节的训练与推理接口，另行记录候选比较和修订成本。

\section{模型能力与使用约束}
\label{sec:vision-models-constraints}

一个模型支持的输入和输出，是开始选型的依据。实际使用还需确认信息经过编码、压缩和生成后，是否保留了当前任务所需的内容，以及达到这一结果用了多少资源。下面从信息、训练依据和成本三个方面回到同一街景与音画请求。

\subsection{信息保留与输出能力}
\label{sec:vision-models-information}

整图缩小会让车牌字符变成几个像素，局部裁剪会删除周围关系，视频抽帧会漏掉短动作，固定查询数会把许多视觉位置汇总成少量表示。投影和共享主干可以重新组织已经保留的信息，却不能保证恢复这些阶段丢弃的证据。

可以固定问题，分别提供整图、局部高分辨率图，以及带原图位置的二者组合。如果只有局部图能正确读取车牌，就应进一步判断错误来自原始分辨率、输入缩放还是后续压缩，而不是立即归为语言推理失败。组合输入还要确认模型知道局部图来自哪里，否则可能把两个视图当成两辆不同的车。

输出也需逐项验收。匹配分数正确只说明当前候选中的选择，语言流畅只说明形式自然，视频画面逼真只说明感知质量的一部分。数量、位置、时间与来源需要对应的检查。精确区域和答案证据见\ref{chap:vision-multimodal-understanding}章，媒体连贯性和控制见\ref{chap:vision-multimodal-generation}章。

\subsection{模态贡献与训练依据}
\label{sec:vision-models-modality-contribution}

模型接受声音，并不说明一次回答确实使用了声音。可以在保持问题不变时移除音频，或换成时间错开的声音，检查结果与正确答案怎样变化。若问题只问车的颜色，移除声音后答案不变是合理的；若问题问敲击发生了几次，错误声音改变了答案，则需要判断模型是否把它与画面进行了适当核验。

同样，替换街景图却保留文字问题，可以检查模型是否依赖常见语言模式猜测。替换条件必须使正确答案或证据关系发生可解释的变化；只观察回答变了，无法证明模型使用了正确模态。多模态评价还需要覆盖缺失模态、冲突模态和训练中未直接出现的组合。

训练依据决定了可以期待哪一层联系。整图配对监督适合建立全局对应，框或时间区间标注提供更细的位置联系，指令示例说明请求和回答怎样组织。它们可以联合使用，但不能互相冒充。某个新组合的成功应在同类样本与清楚的任务协议中复核，不能由一次演示推断任意模态都已对齐。

\subsection{计算资源与使用成本}
\label{sec:vision-models-cost}

多模态成本可以按输入编码、联合上下文计算、输出生成和媒体解码记录。固定图像大小还不够，裁剪数、视频采样、音频长度、历史轮数和输出时长都可能改变实际输入与输出规模。一个只生成短答案的调用，不能与生成数秒音视频的调用仅比较总延迟。

在标准稠密注意力中，$n$个序列位置的成对分数规模为$n^2$，隐藏宽度为$d$时，相关乘加量含$O(n^2d)$项。若有256个文字单元和1024个视觉单元，总长为1280；视觉单元增加到4096时，总长4352，成对位置数约变为原来的11.56倍。这个教学比例只描述该计算项，不能直接当作实际延迟倍率，局部注意力、缓存与分块实现都会改变资源安排。

冻结模块减少需要保存的参数梯度和优化器状态，却不省去推理中的前向计算。若可训练连接器位于冻结语言模型之前，训练还可能需要经过冻结层计算输入梯度。媒体生成则通常反复调用速度或去噪网络，成本取决于步数和每步条件计算；候选重排还增加整条生成路径的重复次数。

共享主干和模型组合也需按实际调用比较。前者可能复用上下文，后者可以独立缓存编码或选择不同规模的专用模块。部署时还需记录编码、缓存和生成各环节的资源消耗，以及等待和输出长度的限制。

\section{本章小结}
\label{sec:vision-models-summary}

街景照片、请求和声音进入同一系统，需要经过有来源与位置约定的编码，再由匹配空间、视觉读取或共享主干建立联系。CLIP与ImageBind说明可比较表示怎样形成；视觉前缀、层间交叉注意力和多层接入说明语言生成怎样利用视觉；潜空间生成与速度预测则为媒体输出提供了不同的计算路径。

这些路径可以在同一模型中共享，也可以通过专用模块组合。语言回答依靠下一词元预测表达判断，图像和音视频输出还需相应潜表示、生成更新与解码器。输入支持范围、训练过的对应关系和实际输出能力应分别确认。

当一次输入不足以回答时，新的裁剪应从原始来源取回像素，而不是只延长解释。CounterCredit通过固定调用前状态的两个对照，把视觉分支的收益与真实返回像素的收益分开，适用范围仍受裁剪参照与评分方式限制。生成中的反思编辑则读取已经形成的图片：UniReason用选择后的文字与图像目标共同训练一轮修订，但修订费用和条件保持还需单独检查。二者都把后续语言计算与新图像联系起来，前者取得真实观测，后者处理创作结果，不能把这两类来源混为证据。

编码和连接中的证据丢失、时间对应不可靠及不同输出目标之间的取舍，仍需要结合具体任务研究。重点是怎样在固定预算中保留必需的局部证据、建立可靠时间对应，并协调不同输出目标；扩大名义容量只是其中一种手段。选好接口后，仍需回到任务检查答案与证据、内容与条件是否一致，这也是后续理解与生成两章的重点。
